\makeatletter
\def\input@path{{./}{../}{../../}{../../../}{../../../../}{preamble/}{../preamble/}{../../preamble/}{../../../preamble/}{../../../../book2026-2027/preamble/}}
\makeatother

\input{preamble_exams.tex}
\input{graphs/highlights.tex}

 % \GTenHighlightsfalse % Uncomment to disable instructional highlights.
 
\input{graphs/specific_highlights.tex}
\CtwoPracticeIntervalPalette
\newcommand{\CtwoSignificance}[1]{\CtwoConfidence{#1}}

% Document-local concise rendering of the existing registry key.
\expandafter\renewcommand\csname practicehighlightlegend@G11-T2-C2\endcsname{%
\begingroup
\renewcommand{\arraystretch}{1.12}
\noindent\begin{tabularx}{\linewidth}{@{}>{\raggedright\arraybackslash}p{.23\linewidth}>{\raggedright\arraybackslash}p{.29\linewidth}>{\raggedright\arraybackslash}X@{}}
\toprule
\textbf{Symbol or notation} & \textbf{Concept} & \textbf{Brief meaning} \\
\midrule
$\CtwoMean{\mu}$ & \CtwoMean{Population mean} & Fixed mean to estimate. \\
$\CtwoPoint{\bar{x}}$ & \CtwoPoint{Point estimate} & Sample mean; interval centre. \\
$\CtwoSize{n}$ & \CtwoSize{Sample size} & Number of sampled units. \\
$\CtwoSD{\sigma}$ & \CtwoSD{Population SD} & Spread of individual observations. \\
$\CtwoSE{\mathrm{SE}}$ & \CtwoSE{Standard error} & Spread of sample means. \\
$\CtwoConfidence{C}$; $\CtwoSignificance{\alpha}$, $\CtwoSignificance{\alpha}/2$ & \CtwoConfidence{Confidence level}; \CtwoSignificance{significance level} & Long-run coverage; total and per-tail areas outside. \\
$\CtwoCritical{z^*}$ & \CtwoCritical{Critical value} & Standard normal multiplier. \\
$\CtwoMoE{\mathrm{MoE}}$ & \CtwoMoE{Margin of error} & Distance from centre to a bound. \\
$\CtwoEndpoint{L}$, $\CtwoEndpoint{U}$ & \CtwoEndpoint{Endpoints} & Lower and upper interval bounds. \\
$\CtwoEndpoint{U}-\CtwoEndpoint{L}$ & \CtwoWidth{Full width} & Distance between bounds. \\
$\CtwoCoverage{\text{rate}}$ & \CtwoCoverage{Observed coverage} & Fraction covering in a finite run. \\
\bottomrule
\end{tabularx}\par
\endgroup}


\ifGTenHighlights
    \colorlet{CtwoCentralArea}{CtwoConfidenceColor!55}
    \colorlet{CtwoCoverLine}{GTenHLOutcome!70!black}
\else
    \colorlet{CtwoCentralArea}{black!25}
    \colorlet{CtwoCoverLine}{black}
\fi

\usepackage{titling}
\setlength{\droptitle}{-2em}
\pretitle{\par\begingroup\centering\LARGE}
\posttitle{\par\endgroup}
\preauthor{}\postauthor{}
\predate{}\postdate{}
\maketitlehookd{\vspace{-1em}}
\linespread{1.12}
\setlength{\parskip}{4pt}
\newgeometry{margin=0.65in,top=0.55in,bottom=0.65in}
\raggedbottom

\title{Grade 11 Term 2 Home Practice -- C2\\\large Point estimates and confidence levels\\[0.25em]
    \normalsize Global Economics -- 2026--2027}
\author{}
\date{}

\newcommand{\InformationTask}{\textbf{Task 1 -- Identify and organize the information; calculate SE.} Make a table with exactly four columns: \textbf{Symbol}, \textbf{Concept}, \textbf{Type}, and \textbf{Meaning in context}. Give each symbol and concept its own row. Use these exact Type labels: \textbf{Sample statistics} for sample mean $\bar{x}$, sample standard deviation $s$ (only when provided or relevant), sample size $n$, and standard error $\mathrm{SE}$; \textbf{Population parameters} for population mean $\mu$ and population standard deviation $\sigma$; \textbf{Procedure Settings} for confidence level $C$, significance level $\alpha$, tail area $\alpha/2$, and critical value $z^*$. Follow this organizational convention for $n$ and $\mathrm{SE}$. Keep Symbol to notation, Concept to the concept name, and Type to its classification. In Meaning in context, identify the population and unknown mean with units, explain $\bar{x}$ as its point estimate, record known $\sigma$, and clearly label shared or differing values for each team or confidence level. Give $C$ as a decimal, $\alpha=1-C$, $\alpha/2$, and $z^*$ in separate rows. Calculate $\mathrm{SE}=\sigma/\sqrt{n}$ with substitution, result, and units in its row.}
\newcommand{\IntervalTask}{\textbf{Task 2 -- Calculate MoE and state the interval.} For each procedure, calculate $\mathrm{MoE}=z^*\times\mathrm{SE}$ and $[\bar{x}-\mathrm{MoE},\bar{x}+\mathrm{MoE}]$. Show substitutions, endpoints, and units; briefly interpret each interval in context. Compare widths when there is more than one interval.}
\newcommand{\CoverageTask}{\textbf{Task 3 -- Calculate observed coverage and answer the multiple-choice question.} Extract the covering count and total from the simulation evidence. Calculate $(\text{covering count}/\text{total})\times100\%$ separately for each procedure, then select one option and explain its interpretation of confidence.}

% Task 1 has its own structure; retain EstimationTable for the other tasks.
\newcommand{\InformationTable}[1]{%
    \begingroup
    \small
    \setlength{\tabcolsep}{4pt}
    \renewcommand{\arraystretch}{1.16}
    \noindent
    \begin{tabularx}{\linewidth}{@{}>{\raggedright\arraybackslash}p{0.07\linewidth}>{\raggedright\arraybackslash}p{0.18\linewidth}>{\raggedright\arraybackslash}p{0.19\linewidth}>{\raggedright\arraybackslash}X@{}}
        \toprule
        \textbf{Symbol} & \textbf{Concept} & \textbf{Type} & \textbf{Meaning in context} \\
        \midrule
        #1
        \bottomrule
    \end{tabularx}\par
    \vspace{2mm}
    \endgroup
}

\newcommand{\EstimationTable}[1]{%
    \noindent
    \begin{tabularx}{\linewidth}{@{}>{\raggedright\arraybackslash}p{0.17\linewidth}>{\raggedright\arraybackslash}p{0.49\linewidth}>{\raggedright\arraybackslash}X@{}}
        \toprule
        \textbf{Quantity} & \textbf{Calculation and result} & \textbf{Meaning in context} \\
        \midrule
        #1
        \bottomrule
    \end{tabularx}\par
    \vspace{2mm}
}

\newlist{answeroptions}{enumerate}{1}
\setlist[answeroptions]{label=\Alph*., leftmargin=4em, topsep=0pt, itemsep=2pt, parsep=0pt}

\begin{document}
\maketitle

\setcounter{tocdepth}{1}
\setcounter{secnumdepth}{2}
\phantomsection
\label{toc}
\tableofcontents

\section*{Evaluated criteria}
\addcontentsline{toc}{section}{Evaluated criteria}
\begin{description}
    \item[\textbf{C2}] Compares point and interval estimation and interprets confidence levels through their long-run coverage meaning, distinguishing this from the probability that a particular calculated interval contains the parameter.
\end{description}

\textbf{How to practise.} Try each task before opening its blue solution link. Use the arrows to return to the contents, problem list, or problem. Use the same three tasks throughout: (1) organize the information in a table and calculate \CtwoSE{SE}; (2) calculate \CtwoMoE{MoE}, state and interpret each interval; (3) calculate \CtwoCoverage{observed coverage} from the evidence and answer one multiple-choice question.

\clearpage
\section{Introduction}
\label{sec:introduction}

\noindent\begin{minipage}{\linewidth}
\subsection{Key terms and highlight guide}
\label{subsec:c2-key-terms}
\phantomsection\label{tab:highlight-legend}
\PracticeHighlightLegend{G11-T2-C2}
\end{minipage}\par

\phantomsection\label{toc:c2-introduction}
\begingroup
% Relabel only the shared index's bold heading, within this local group.
\let\CtwoContentsBold\textbf
\renewcommand{\textbf}[1]{\ifstrequal{#1}{Problems:}{\CtwoContentsBold{Introduction contents:}}{\CtwoContentsBold{#1}}}
\SectionProblemsTableOfContents{
    \SectionProblemLink{subsec:c2-key-terms}{Key terms and highlight guide} \\
    \SectionProblemLink{subsec:c2-cycling-example}{Motivating example 1: a cycling estimate} \\
    \SectionProblemLink{subsec:c2-coverage-example}{Motivating example 2: repeated samples} \\
    \SectionProblemLink{subsec:c2-comprehensive}{Worked example: investigate the cycling interval}
}
\endgroup
\par\medskip
\noindent\small Introduction contents. The arrows below return to the main contents,
this list, and (within the worked example) its opening.\normalsize
\par\vspace{3mm}\secInicio[][toc:c2-introduction]
\clearpage


\noindent\begin{minipage}{\linewidth}
\subsection{Example 1: from cycling trips to an interval}
\label{subsec:c2-cycling-example}
A planner wants the \CtwoMean{population mean} $\CtwoMean{\mu}$ of weekly cycling-trip times to the market. An independent random sample of $\CtwoSize{n}=\CtwoSize{49}$ trips gives the \CtwoPoint{point estimate} $\CtwoPoint{\bar{x}}=\CtwoPoint{40\,\text{minutes}}$. Assume normal trip times with known \CtwoSD{population SD} $\CtwoSD{\sigma}=\CtwoSD{14\,\text{minutes}}$. Then the \CtwoSE{standard error} is
\[
\CtwoSE{\mathrm{SE}}=\frac{\CtwoSD{14}}{\sqrt{\CtwoSize{49}}}=\CtwoSE{2\,\text{minutes}}.
\]
For \CtwoConfidence{95\%} confidence, $\CtwoConfidence{C}=\CtwoConfidence{0.95}$, $\CtwoSignificance{\alpha}=1-\CtwoConfidence{0.95}=\CtwoSignificance{0.05}$ and each tail has area $\CtwoSignificance{\alpha}/2=\CtwoSignificance{0.025}$. The \CtwoCritical{critical value} $\CtwoCritical{z^*}=\CtwoCritical{1.960}$ gives \CtwoMoE{margin of error} $\CtwoMoE{\mathrm{MoE}}=\CtwoCritical{1.960}\times\CtwoSE{2}=\CtwoMoE{3.92\,\text{minutes}}$. Thus $\CtwoPoint{40}\pm\CtwoMoE{3.92}$ has \CtwoEndpoint{endpoints} $\CtwoEndpoint{L}=\CtwoEndpoint{36.08}$ and $\CtwoEndpoint{U}=\CtwoEndpoint{43.92}$ minutes, and \CtwoWidth{full width} $2(\CtwoMoE{3.92})=\CtwoWidth{7.84\,\text{minutes}}$.

The interval gives plausible values for the fixed mean under the model. Since it includes 42 and larger values, it does not support concluding that the mean is under 42 minutes.
\par\vspace{3mm}\secInicio[][toc:c2-introduction]
\end{minipage}\par

\noindent\begin{minipage}{\linewidth}
\subsection{Example 2: repeating the cycling procedure}
\label{subsec:c2-coverage-example}
A separate, hypothetical normal population has known fixed mean $\CtwoMean{\mu}=\CtwoMean{41\,\text{minutes}}$ and the same $\CtwoSD{\sigma}=\CtwoSD{14\,\text{minutes}}$. Simulate 50 fresh independent samples of $\CtwoSize{n}=\CtwoSize{49}$ trips and apply Example 1's \CtwoConfidence{95\%} procedure to each. Each new $\CtwoPoint{\bar{x}}$ gives new endpoints; the population mean stays fixed. If 3 intervals miss 41, then 47 cover, so the \CtwoCoverage{observed coverage} is $(47/50)\times100\%=\CtwoCoverage{94\%}$.

This finite-run rate does not change $\CtwoConfidence{C}=\CtwoConfidence{0.95}$ or $\CtwoSignificance{\alpha}=\CtwoSignificance{0.05}$. Under the model, about \CtwoConfidence{95\%} of intervals cover over many repetitions. Each realized interval either contains the mean or does not; confidence does not give the percentage of individual trip times inside it.

\textbf{Quick check.} Which can change in a new run: the sample means, the observed coverage, or the chosen confidence level? The first two can change; the level stays fixed unless the procedure is changed.
\par\vspace{3mm}\secInicio[subsec:c2-cycling-example][toc:c2-introduction]
\end{minipage}\par

\clearpage

\subsection{Comprehensive worked example: investigate the cycling interval}
\label{subsec:c2-comprehensive}
We now unpack Example 1, then change one ingredient at a time. All six graphical
investigations belong to this one worked example.

\subsubsection*{Step 1. Identify what is fixed and what is sampled}
The \textbf{population} is all cycling trips to the market in the planner's target
week. Its \CtwoMean{population mean} $\CtwoMean{\mu}$ is one fixed, unknown number
of minutes. The \textbf{sample} is the \CtwoSize{49 trips} selected at random.
Their \CtwoPoint{sample mean} $\CtwoPoint{\bar{x}}=\CtwoPoint{40\,\text{minutes}}$
is a \CtwoPoint{point estimate}: one number used to estimate that unknown mean.
A different sample would usually give a different estimate.

The known \CtwoSD{population standard deviation}
$\CtwoSD{\sigma}=\CtwoSD{14\,\text{minutes}}$ measures variation among
\emph{individual trip times}. Averaging independent observations reduces variation:
the \CtwoSE{standard error} describes the spread of \emph{sample means} over
repeated samples of the same size, rather than the spread of the individual trips.
\[
\CtwoSE{\mathrm{SE}}=\frac{\CtwoSD{\sigma}}{\sqrt{\CtwoSize{n}}}
=\frac{\CtwoSD{14}}{\sqrt{\CtwoSize{49}}}
=\frac{\CtwoSD{14}}{7}=\CtwoSE{2\,\text{minutes}}.
\]
Thus means of 49 trips vary less than single trips: \CtwoSE{2 minutes} versus
\CtwoSD{14 minutes}. SE is a standard deviation, not the observed error
$\bar{x}-\mu$, which we cannot calculate while $\mu$ is unknown.

\textbf{Why a two-sided $z$ interval?} Assume a random, representative sample,
independent trip times, a genuinely known population SD, and normal trip times.
These assumptions make the sample mean normal with standard deviation
$\sigma/\sqrt n$. For a non-normal population, a sufficiently large sample can
make this an approximation, provided severe skew or outliers do not undermine it;
$n=49$ alone is not a guarantee. When sampling without replacement, the population
should be much larger than the sample (for example, at least ten times as large).
Two-sided means that uncertainty is allowed both below and above the estimate.

\subsubsection*{Step 2. Choose confidence and locate the two normal tails (E)}
Choose \CtwoConfidence{confidence level} $\CtwoConfidence{C}=\CtwoConfidence{0.95}$.
The remaining standard normal area is
$\CtwoSignificance{\alpha}=1-\CtwoConfidence{0.95}=\CtwoSignificance{0.05}$,
split equally into $\CtwoSignificance{\alpha}/2=\CtwoSignificance{0.025}$ per tail.
The positive \CtwoCritical{critical value} has left area $1-0.025=0.975$:
\[
\CtwoCritical{z^*}=\Phi^{-1}(1-\CtwoSignificance{\alpha}/2)
=\Phi^{-1}(0.975)\approx\CtwoCritical{1.960}.
\]
Here $\Phi(z)$ is standard normal area to the left of $z$; its inverse finds the cutoff.

\begin{center}
\begin{tikzpicture}
\begin{axis}[width=.86\linewidth,height=4.6cm,axis x line=bottom,axis y line=left,
 xmin=-4,xmax=4,ymin=0,ymax=.45,xlabel={Standardized sample mean $z=(\bar{x}-\mu)/\mathrm{SE}$ (no units)},
 ylabel={Density},ytick={0,.2,.4},xtick={-1.96,0,1.96},
 xticklabels={$-1.960$,$0$,$1.960$},tick label style={font=\small},
 samples=121,declare function={nd(\x)=exp(-\x^2/2)/sqrt(2*pi);}]
\addplot[draw=none,fill=black!12,domain=-4:-1.96] {nd(x)} \closedcycle;
\addplot[draw=none,fill=CtwoCentralArea,domain=-1.96:1.96] {nd(x)} \closedcycle;
\addplot[draw=none,fill=black!12,domain=1.96:4] {nd(x)} \closedcycle;
\addplot[black,thick,domain=-4:4] {nd(x)};
\addplot[dashed] coordinates {(-1.96,0) (-1.96,.05844)};
\addplot[dashed] coordinates {(1.96,0) (1.96,.05844)};
\node at (axis cs:0,.18) {\small\CtwoConfidence{Central 95\%}};
\node[font=\small] (lt) at (axis cs:-2.75,.29) {\CtwoSignificance{2.5\%}};
\node[font=\small] (rt) at (axis cs:2.75,.29) {\CtwoSignificance{2.5\%}};
\draw[->] (lt.south) -- (axis cs:-2.3,.015);
\draw[->] (rt.south) -- (axis cs:2.3,.015);
\end{axis}
\end{tikzpicture}
\par\small\textbf{E.} Areas describe the standardized \emph{sample mean} under the model,
not the distribution of individual cycling times.
\end{center}
\par\vspace{3mm}\secInicio[subsec:c2-comprehensive][toc:c2-introduction]

\clearpage
\subsubsection*{Step 3. Build the interval and interpret it}
The normal cutoffs tell us how many standard errors to allow on each side.
Multiplying by SE converts that unitless multiplier into minutes:
\[
\CtwoMoE{\mathrm{MoE}}=\CtwoCritical{z^*}\times\CtwoSE{\mathrm{SE}}
=\CtwoCritical{1.960}\times\CtwoSE{2}=\CtwoMoE{3.92\,\text{minutes}}.
\]
Centre this allowance on the observed mean, not on the unknown population mean:
\begin{align*}
[\CtwoEndpoint{L},\CtwoEndpoint{U}]
&=[\CtwoPoint{\bar{x}}-\CtwoMoE{\mathrm{MoE}},\CtwoPoint{\bar{x}}+\CtwoMoE{\mathrm{MoE}}]\\
&=[\CtwoPoint{40}-\CtwoMoE{3.92},\CtwoPoint{40}+\CtwoMoE{3.92}]
=[\CtwoEndpoint{36.08},\CtwoEndpoint{43.92}]\ \text{minutes}.\\
\text{Centre}&=(36.08+43.92)/2=\CtwoPoint{40\,\text{minutes}},\\
\CtwoWidth{\text{Full width}}&=\CtwoEndpoint{43.92}-\CtwoEndpoint{36.08}
=2\CtwoMoE{(3.92)}=\CtwoWidth{7.84\,\text{minutes}}.
\end{align*}
\textbf{In context:} we are \CtwoConfidence{95\% confident} that the population's
mean cycling-trip time is between \CtwoEndpoint{36.08} and \CtwoEndpoint{43.92 minutes}.
This is a range of plausible values for the mean under the assumptions.
The confidence belongs to the method: in many repetitions, about
\CtwoConfidence{95\%} of the intervals would contain the fixed mean.

For this particular interval, the fixed mean either lies inside or lies outside.
We do \emph{not} assign a 95\% probability to the fixed mean being inside this
already calculated interval. Nor does this interval contain 95\% of individual trips.

\subsubsection*{Step 4. Compare confidence levels for the same sample (F)}
Keep $\CtwoPoint{\bar{x}}=\CtwoPoint{40}$ and
$\CtwoSE{\mathrm{SE}}=\CtwoSE{2}$ minutes. Only the chosen confidence level changes.
Using critical values rounded to three decimals and endpoints to two decimals:
\begin{center}
\begin{tabular}{@{}cccc@{}}
\toprule
$\CtwoConfidence{C}$ & $\CtwoCritical{z^*}$ & $\CtwoMoE{\mathrm{MoE}}=\CtwoCritical{z^*}(\CtwoSE{2})$ (min) & Interval (min)\\
\midrule
\CtwoConfidence{90\%} & \CtwoCritical{1.645} & \CtwoMoE{3.290} & $[\CtwoEndpoint{36.71},\CtwoEndpoint{43.29}]$\\
\CtwoConfidence{95\%} & \CtwoCritical{1.960} & \CtwoMoE{3.920} & $[\CtwoEndpoint{36.08},\CtwoEndpoint{43.92}]$\\
\CtwoConfidence{99\%} & \CtwoCritical{2.576} & \CtwoMoE{5.152} & $[\CtwoEndpoint{34.85},\CtwoEndpoint{45.15}]$\\
\bottomrule
\end{tabular}
\end{center}
\begin{center}
\begin{tikzpicture}
\begin{axis}[width=.86\linewidth,height=4cm,xmin=33,xmax=47,ymin=.5,ymax=3.6,
 xlabel={Population mean trip time being estimated (minutes)},ytick={1,2,3},
 yticklabels={99\%,95\%,90\%},ylabel={Confidence},xtick={34,36,38,40,42,44,46},
 tick label style={font=\small},clip=false]
\addplot[densely dotted] coordinates {(40,.6) (40,3.5)};
\addplot[very thick,mark=|] coordinates {(36.71,3) (43.29,3)};
\addplot[very thick,dashed,mark=|] coordinates {(36.08,2) (43.92,2)};
\addplot[very thick,dashdotted,mark=|] coordinates {(34.848,1) (45.152,1)};
\addplot[only marks,mark=*,mark size=2pt] coordinates {(40,1) (40,2) (40,3)};
\node[font=\small,anchor=south] at (axis cs:40,3.5) {Same centre: \CtwoPoint{40 min}};
\end{axis}
\end{tikzpicture}
\par\small\textbf{F.} Endpoints are labelled in the table; dots mark the common sample mean.
\end{center}
The intervals are \textbf{nested}: a larger critical value extends both ends while
keeping the same centre. These are three procedures applied to \emph{one sample},
not three different samples. Different samples need not produce nested intervals.
\par\vspace{3mm}\secInicio[subsec:c2-comprehensive][toc:c2-introduction]

\clearpage
\subsubsection*{Step 5. See what confidence means across repeated samples (A)}
For checking coverage only, switch to a hypothetical normal cycling population
with \emph{known} $\CtwoMean{\mu}=\CtwoMean{41\,\text{minutes}}$.
This is a simulation setting, separate from the planner's problem where
$\CtwoMean{\mu}$ is unknown. Keep $\CtwoSD{\sigma}=\CtwoSD{14}$,
$\CtwoSize{n}=\CtwoSize{49}$ and $\CtwoConfidence{C}=\CtwoConfidence{0.95}$.

Each row uses an explicitly specified illustrative sample mean, shown on the right.
Every interval is calculated by the same rule:
$[L_i,U_i]=[\bar{x}_i-3.92,\bar{x}_i+3.92]$ minutes.
For example, sample 1 gives $[38.20-3.92,38.20+3.92]=[34.28,42.12]$;
sample 6 gives $[36.50-3.92,36.50+3.92]=[32.58,40.42]$ and misses 41.

\begin{center}
\begin{tikzpicture}
\begin{axis}[width=.90\linewidth,height=10.6cm,xmin=32,xmax=53,ymin=.3,ymax=21.3,
 y dir=reverse,ytick={1,...,20},ylabel={Sample number},
 xtick={32,36,40,44,48},xlabel={Interval endpoints and sample mean (minutes)},
 tick label style={font=\small},clip=false]
\addplot[black,densely dotted,thick] coordinates {(41,.5) (41,20.5)};
\node[font=\small,anchor=south] at (axis cs:41,.5) {\CtwoMean{$\mu=41$}};
\node[font=\small,anchor=south] at (axis cs:51,.5) {$\bar{x}_i$ (min)};
\addplot[CtwoCoverLine,thick,solid,mark=|] coordinates {(34.28,1) (42.12,1)};
\addplot[CtwoCoverLine,only marks,mark=*,mark size=2pt] coordinates {(38.20,1)};
\node[font=\small] at (axis cs:51,1) {38.20};
\addplot[CtwoCoverLine,thick,solid,mark=|] coordinates {(38.18,2) (46.02,2)};
\addplot[CtwoCoverLine,only marks,mark=*,mark size=2pt] coordinates {(42.10,2)};
\node[font=\small] at (axis cs:51,2) {42.10};
\addplot[CtwoCoverLine,thick,solid,mark=|] coordinates {(36.58,3) (44.42,3)};
\addplot[CtwoCoverLine,only marks,mark=*,mark size=2pt] coordinates {(40.50,3)};
\node[font=\small] at (axis cs:51,3) {40.50};
\addplot[CtwoCoverLine,thick,solid,mark=|] coordinates {(37.88,4) (45.72,4)};
\addplot[CtwoCoverLine,only marks,mark=*,mark size=2pt] coordinates {(41.80,4)};
\node[font=\small] at (axis cs:51,4) {41.80};
\addplot[CtwoCoverLine,thick,solid,mark=|] coordinates {(35.68,5) (43.52,5)};
\addplot[CtwoCoverLine,only marks,mark=*,mark size=2pt] coordinates {(39.60,5)};
\node[font=\small] at (axis cs:51,5) {39.60};
\addplot[black,thick,dashed,mark=|] coordinates {(32.58,6) (40.42,6)};
\addplot[only marks,mark=x,mark size=2pt] coordinates {(36.50,6)};
\node[font=\small] at (axis cs:51,6) {36.50};
\addplot[CtwoCoverLine,thick,solid,mark=|] coordinates {(39.28,7) (47.12,7)};
\addplot[CtwoCoverLine,only marks,mark=*,mark size=2pt] coordinates {(43.20,7)};
\node[font=\small] at (axis cs:51,7) {43.20};
\addplot[CtwoCoverLine,thick,solid,mark=|] coordinates {(36.18,8) (44.02,8)};
\addplot[CtwoCoverLine,only marks,mark=*,mark size=2pt] coordinates {(40.10,8)};
\node[font=\small] at (axis cs:51,8) {40.10};
\addplot[CtwoCoverLine,thick,solid,mark=|] coordinates {(37.18,9) (45.02,9)};
\addplot[CtwoCoverLine,only marks,mark=*,mark size=2pt] coordinates {(41.10,9)};
\node[font=\small] at (axis cs:51,9) {41.10};
\addplot[CtwoCoverLine,thick,solid,mark=|] coordinates {(38.78,10) (46.62,10)};
\addplot[CtwoCoverLine,only marks,mark=*,mark size=2pt] coordinates {(42.70,10)};
\node[font=\small] at (axis cs:51,10) {42.70};
\addplot[CtwoCoverLine,thick,solid,mark=|] coordinates {(34.98,11) (42.82,11)};
\addplot[CtwoCoverLine,only marks,mark=*,mark size=2pt] coordinates {(38.90,11)};
\node[font=\small] at (axis cs:51,11) {38.90};
\addplot[CtwoCoverLine,thick,solid,mark=|] coordinates {(36.88,12) (44.72,12)};
\addplot[CtwoCoverLine,only marks,mark=*,mark size=2pt] coordinates {(40.80,12)};
\node[font=\small] at (axis cs:51,12) {40.80};
\addplot[CtwoCoverLine,thick,solid,mark=|] coordinates {(39.88,13) (47.72,13)};
\addplot[CtwoCoverLine,only marks,mark=*,mark size=2pt] coordinates {(43.80,13)};
\node[font=\small] at (axis cs:51,13) {43.80};
\addplot[CtwoCoverLine,thick,solid,mark=|] coordinates {(35.18,14) (43.02,14)};
\addplot[CtwoCoverLine,only marks,mark=*,mark size=2pt] coordinates {(39.10,14)};
\node[font=\small] at (axis cs:51,14) {39.10};
\addplot[CtwoCoverLine,thick,solid,mark=|] coordinates {(38.48,15) (46.32,15)};
\addplot[CtwoCoverLine,only marks,mark=*,mark size=2pt] coordinates {(42.40,15)};
\node[font=\small] at (axis cs:51,15) {42.40};
\addplot[CtwoCoverLine,thick,solid,mark=|] coordinates {(36.38,16) (44.22,16)};
\addplot[CtwoCoverLine,only marks,mark=*,mark size=2pt] coordinates {(40.30,16)};
\node[font=\small] at (axis cs:51,16) {40.30};
\addplot[black,thick,dashed,mark=|] coordinates {(41.38,17) (49.22,17)};
\addplot[only marks,mark=x,mark size=2pt] coordinates {(45.30,17)};
\node[font=\small] at (axis cs:51,17) {45.30};
\addplot[CtwoCoverLine,thick,solid,mark=|] coordinates {(37.58,18) (45.42,18)};
\addplot[CtwoCoverLine,only marks,mark=*,mark size=2pt] coordinates {(41.50,18)};
\node[font=\small] at (axis cs:51,18) {41.50};
\addplot[CtwoCoverLine,thick,solid,mark=|] coordinates {(33.48,19) (41.32,19)};
\addplot[CtwoCoverLine,only marks,mark=*,mark size=2pt] coordinates {(37.40,19)};
\node[font=\small] at (axis cs:51,19) {37.40};
\addplot[CtwoCoverLine,thick,solid,mark=|] coordinates {(40.68,20) (48.52,20)};
\addplot[CtwoCoverLine,only marks,mark=*,mark size=2pt] coordinates {(44.60,20)};
\node[font=\small] at (axis cs:51,20) {44.60};
\end{axis}
\end{tikzpicture}
\par\small\textbf{A. Illustrative repeated-sampling data (20 samples).}
Solid intervals with dots cover; dashed intervals with crosses miss.
The dotted vertical line is the fixed simulation mean.
\end{center}
Only samples 6 and 17 miss: $\CtwoCoverage{18}$ of $20$ intervals meet or cross 41.
Thus \CtwoCoverage{observed coverage} is
$(18/20)\times100\%=\CtwoCoverage{90\%}$, whereas the
\CtwoConfidence{chosen confidence level} remains \CtwoConfidence{95\%}.
This small illustrative collection need not have exactly 19 covering intervals:
counts fluctuate from run to run, and these displayed means are illustrative,
not a claim about a measured dataset.

The population mean stays fixed; the means and endpoints vary across samples.
Changing only $\bar{x}$ shifts the whole interval by the same amount, without
changing its \CtwoWidth{7.84-minute width}. Equal width does not guarantee coverage.
\par\vspace{3mm}\secInicio[subsec:c2-comprehensive][toc:c2-introduction]

\clearpage
\subsubsection*{Step 6. Change spread or sample size, one at a time (B and D)}
\textbf{B. Population spread.} Hold $\CtwoSize{n}=\CtwoSize{49}$ and
$\CtwoConfidence{C}=\CtwoConfidence{95\%}$ fixed. Then
$\CtwoMoE{\mathrm{MoE}}=\CtwoCritical{1.960}\CtwoSD{\sigma}/7$.
This compares hypothetical populations or scenarios: $\sigma$ is fixed within
one population, rather than changing when we take another sample.
\begin{center}
\begin{tikzpicture}
\begin{axis}[width=.85\linewidth,height=5cm,xmin=0,xmax=30,ymin=0,ymax=9,
 xlabel={Population standard deviation $\sigma$ (minutes)},ylabel={MoE (minutes)},
 xtick={0,7,14,21,28},ytick={0,2,4,6,8},grid=major,tick label style={font=\small}]
\addplot[black,thick,domain=.1:30] {1.960*x/7};
\addplot[only marks,mark=*] coordinates {(14,3.92) (28,7.84)};
\node[font=\small,anchor=north west,fill=white] at (axis cs:14,3.92) {$(14,\CtwoMoE{3.92})$};
\node[font=\small,anchor=south east,fill=white] at (axis cs:28,7.84) {$(28,7.84)$};
\end{axis}
\end{tikzpicture}
\par\small\textbf{B.} A straight-line relationship for positive $\sigma$; zero is an axis reference.
\end{center}
Each additional minute of population SD adds $1.960/7=0.28$ minutes to MoE.
Doubling $\sigma$ from 14 to 28 doubles SE and hence MoE from
\CtwoMoE{3.92 to 7.84 minutes}. More variable individual times make the mean harder to estimate precisely.

\textbf{D. Sample size.} Now hold $\CtwoSD{\sigma}=\CtwoSD{14}$ and
$\CtwoConfidence{C}=\CtwoConfidence{95\%}$ fixed:
$\CtwoMoE{\mathrm{MoE}}=\CtwoCritical{1.960}(\CtwoSD{14})/\sqrt{\CtwoSize{n}}$.
\begin{center}
\begin{tikzpicture}
\begin{axis}[width=.85\linewidth,height=5cm,xmin=0,xmax=225,ymin=0,ymax=15,
 xlabel={Sample size $n$ (trips)},ylabel={MoE (minutes)},xtick={0,49,100,150,196,225},
 ytick={0,2,4,6,8,10,12,14},grid=major,tick label style={font=\small}]
\addplot[black,thick,domain=4:225,samples=150] {1.960*14/sqrt(x)};
\addplot[only marks,mark=*] coordinates {(49,3.92) (196,1.96)};
\node[font=\small,anchor=south west,fill=white] at (axis cs:49,3.92) {$(49,\CtwoMoE{3.92})$};
\node[font=\small,anchor=south east,fill=white] at (axis cs:196,1.96) {$(196,\CtwoMoE{1.96})$};
\end{axis}
\end{tikzpicture}
\par\small\textbf{D.} The curve shows the trend for $4\leq n\leq225$; actual sample sizes are whole numbers.
\end{center}
Since $\sqrt{196}=14=2\sqrt{49}$, quadrupling the sample size halves SE
from 2 to 1 minute and halves MoE from \CtwoMoE{3.92 to 1.96 minutes}.
The curve flattens: the same extra number of trips gives smaller reductions
as $n$ grows. These are \textbf{diminishing returns}, not a linear decrease.
\par\vspace{3mm}\secInicio[subsec:c2-comprehensive][toc:c2-introduction]

\clearpage
\subsubsection*{Step 7. Change confidence while holding SE fixed (C)}
Keep $\CtwoSD{\sigma}=\CtwoSD{14}$ and $\CtwoSize{n}=\CtwoSize{49}$,
so $\CtwoSE{\mathrm{SE}}=\CtwoSE{2}$ minutes. The critical value changes with confidence:
\[
\CtwoMoE{\mathrm{MoE}}(C)=\Phi^{-1}\!\left(\frac{1+C}{2}\right)\times\CtwoSE{2},
\qquad 0<C<1.
\]
The formula uses $C$ as a decimal; the graph labels confidence as a percentage.
For instance, 95\% means $C=0.95$, with left area $(1+0.95)/2=0.975$.
\begin{center}
\begin{tikzpicture}
\begin{axis}[width=.88\linewidth,height=6.4cm,xmin=50,xmax=100,ymin=0,ymax=7,
 xlabel={Confidence level (\%)},ylabel={MoE (minutes)},
 xtick={50,60,70,80,90,95,99},ytick={0,1,2,3,4,5,6,7},grid=major,
 tick label style={font=\small}]
% Normal inverse CDF evaluated numerically, not a straight-line approximation.
\addplot[black,thick] coordinates {(50.00,1.34897950) (50.25,1.35685714) (50.50,1.36475588) (50.75,1.37267597) (51.00,1.38061765) (51.25,1.38858115) (51.50,1.39656673) (51.75,1.40457464) (52.00,1.41260513) (52.25,1.42065845) (52.50,1.42873488) (52.75,1.43683467) (53.00,1.44495810) (53.25,1.45310544) (53.50,1.46127697) (53.75,1.46947296) (54.00,1.47769370) (54.25,1.48593948) (54.50,1.49421060) (54.75,1.50250736) (55.00,1.51083005) (55.25,1.51917899) (55.50,1.52755449) (55.75,1.53595686) (56.00,1.54438643) (56.25,1.55284352) (56.50,1.56132847) (56.75,1.56984162) (57.00,1.57838331) (57.25,1.58695388) (57.50,1.59555369) (57.75,1.60418311) (58.00,1.61284249) (58.25,1.62153222) (58.50,1.63025267) (58.75,1.63900422) (59.00,1.64778726) (59.25,1.65660220) (59.50,1.66544944) (59.75,1.67432939) (60.00,1.68324247) (60.25,1.69218910) (60.50,1.70116973) (60.75,1.71018479) (61.00,1.71923473) (61.25,1.72832001) (61.50,1.73744109) (61.75,1.74659846) (62.00,1.75579259) (62.25,1.76502398) (62.50,1.77429312) (62.75,1.78360053) (63.00,1.79294673) (63.25,1.80233225) (63.50,1.81175762) (63.75,1.82122341) (64.00,1.83073018) (64.25,1.84027848) (64.50,1.84986892) (64.75,1.85950208) (65.00,1.86917858) (65.25,1.87889903) (65.50,1.88866407) (65.75,1.89847434) (66.00,1.90833051) (66.25,1.91823323) (66.50,1.92818321) (66.75,1.93818115) (67.00,1.94822775) (67.25,1.95832376) (67.50,1.96846992) (67.75,1.97866700) (68.00,1.98891577) (68.25,1.99921703) (68.50,2.00957161) (68.75,2.01998034) (69.00,2.03044407) (69.25,2.04096367) (69.50,2.05154004) (69.75,2.06217410) (70.00,2.07286678) (70.25,2.08361904) (70.50,2.09443186) (70.75,2.10530625) (71.00,2.11624324) (71.25,2.12724388) (71.50,2.13830925) (71.75,2.14944048) (72.00,2.16063868) (72.25,2.17190504) (72.50,2.18324073) (72.75,2.19464701) (73.00,2.20612511) (73.25,2.21767634) (73.50,2.22930203) (73.75,2.24100353) (74.00,2.25278226) (74.25,2.26463964) (74.50,2.27657716) (74.75,2.28859635) (75.00,2.30069876) (75.25,2.31288601) (75.50,2.32515975) (75.75,2.33752169) (76.00,2.34997358) (76.25,2.36251724) (76.50,2.37515453) (76.75,2.38788736) (77.00,2.40071772) (77.25,2.41364764) (77.50,2.42667924) (77.75,2.43981470) (78.00,2.45305624) (78.25,2.46640620) (78.50,2.47986696) (78.75,2.49344100) (79.00,2.50713088) (79.25,2.52093924) (79.50,2.53486883) (79.75,2.54892248) (80.00,2.56310313) (80.25,2.57741382) (80.50,2.59185769) (80.75,2.60643803) (81.00,2.62115822) (81.25,2.63602179) (81.50,2.65103240) (81.75,2.66619384) (82.00,2.68151007) (82.25,2.69698519) (82.50,2.71262349) (82.75,2.72842942) (83.00,2.74440762) (83.25,2.76056293) (83.50,2.77690039) (83.75,2.79342529) (84.00,2.81014312) (84.25,2.82705964) (84.50,2.84418086) (84.75,2.86151310) (85.00,2.87906294) (85.25,2.89683731) (85.50,2.91484348) (85.75,2.93308905) (86.00,2.95158206) (86.25,2.97033091) (86.50,2.98934450) (86.75,3.00863217) (87.00,3.02820378) (87.25,3.04806975) (87.50,3.06824109) (87.75,3.08872945) (88.00,3.10954719) (88.25,3.13070739) (88.50,3.15222395) (88.75,3.17411166) (89.00,3.19638628) (89.25,3.21906458) (89.50,3.24216450) (89.75,3.26570521) (90.00,3.28970725) (90.25,3.31419267) (90.50,3.33918515) (90.75,3.36471023) (91.00,3.39079542) (91.25,3.41747052) (91.50,3.44476778) (91.75,3.47272227) (92.00,3.50137214) (92.25,3.53075908) (92.50,3.56092868) (92.75,3.59193105) (93.00,3.62382135) (93.25,3.65666053) (93.50,3.69051623) (93.75,3.72546373) (94.00,3.76158722) (94.25,3.79898122) (94.50,3.83775245) (94.75,3.87802198) (95.00,3.91992797) (95.25,3.96362911) (95.50,4.00930892) (95.75,4.05718133) (96.00,4.10749782) (96.25,4.16055691) (96.50,4.21671680) (96.75,4.27641268) (97.00,4.34018076) (97.25,4.40869258) (97.50,4.48280546) (97.75,4.56363897) (98.00,4.65269575) (98.25,4.75206168) (98.50,4.86475812) (98.75,4.99541095) (99.00,5.15165861) (99.01,5.15860536) (99.02,5.16561490) (99.03,5.17268848) (99.04,5.17982737) (99.05,5.18703286) (99.06,5.19430632) (99.07,5.20164913) (99.08,5.20906273) (99.09,5.21654860) (99.10,5.22410828) (99.11,5.23174334) (99.12,5.23945542) (99.13,5.24724620) (99.14,5.25511742) (99.15,5.26307089) (99.16,5.27110848) (99.17,5.27923211) (99.18,5.28744378) (99.19,5.29574557) (99.20,5.30413962) (99.21,5.31262815) (99.22,5.32121348) (99.23,5.32989799) (99.24,5.33868418) (99.25,5.34757463) (99.26,5.35657202) (99.27,5.36567914) (99.28,5.37489889) (99.29,5.38423431) (99.30,5.39368852) (99.31,5.40326482) (99.32,5.41296662) (99.33,5.42279750) (99.34,5.43276117) (99.35,5.44286152) (99.36,5.45310263) (99.37,5.46348875) (99.38,5.47402433) (99.39,5.48471404) (99.40,5.49556277) (99.41,5.50657565) (99.42,5.51775807) (99.43,5.52911569) (99.44,5.54065447) (99.45,5.55238067) (99.46,5.56430091) (99.47,5.57642214) (99.48,5.58875174) (99.49,5.60129746) (99.50,5.61406754) (99.51,5.62707068) (99.52,5.64031612) (99.53,5.65381367) (99.54,5.66757374) (99.55,5.68160744) (99.56,5.69592657) (99.57,5.71054378) (99.58,5.72547253) (99.59,5.74072727) (99.60,5.75632348) (99.61,5.77227779) (99.62,5.78860811) (99.63,5.80533371) (99.64,5.82247545) (99.65,5.84005589) (99.66,5.85809950) (99.67,5.87663290) (99.68,5.89568510) (99.69,5.91528783) (99.70,5.93547585) (99.71,5.95628737) (99.72,5.97776453) (99.73,5.99995399) (99.74,6.02290752) (99.75,6.04668288) (99.76,6.07134473) (99.77,6.09696581) (99.78,6.12362830) (99.79,6.15142564) (99.80,6.18046461) (99.81,6.21086805) (99.82,6.24277830) (99.83,6.27636154) (99.84,6.31181352) (99.85,6.34936705) (99.86,6.38930211) (99.87,6.43195952) (99.88,6.47776024) (99.89,6.52723274) (99.90,6.58105346)};
\addplot[only marks,mark=*] coordinates {(90,3.28970725) (95,3.91992797) (99,5.15165861)};
\node[font=\small,anchor=south east,fill=white] at (axis cs:90,3.289707) {90\%: 3.29};
\node[font=\small,anchor=south east,fill=white] at (axis cs:95,3.919928) {95\%: 3.92};
\node[font=\small,anchor=south east,fill=white] at (axis cs:99,5.151659) {99\%: 5.15};
\end{axis}
\end{tikzpicture}
\par\small\textbf{C.} Inverse-normal values computed on $50\%\leq C\leq99.9\%$.
There is no plotted value at 100\%; point labels are rounded to two decimals.
\end{center}
The critical values at 90\%, 95\%, and 99\% are approximately
\CtwoCritical{1.645}, \CtwoCritical{1.960}, and \CtwoCritical{2.576}.
The curve rises nonlinearly because equal extra areas near the thin normal tails
require increasingly large changes in the cutoff. Greater confidence requires
a larger multiplier, a larger MoE, and a wider interval for the same sample.
At 100\%, no finite normal cutoff encloses all the area.

\textbf{Read the investigations together.} A smaller population spread, a larger
sample, or a lower confidence level can narrow an interval. A \CtwoWidth{narrower
interval} means greater precision under the model; it does not guarantee coverage.
Increasing the sample size also does not eliminate sampling bias: consistently
selecting unusually fast cyclists can give a precise estimate of the wrong group.

\textbf{Final check.} Changing only the sample mean moves the interval; changing
SE or confidence changes its width. The fixed population mean does neither.
The selected confidence level describes long-run coverage; the observed coverage
counts successes in one finite collection of intervals.
\par\vspace{3mm}\secInicio[subsec:c2-comprehensive][toc:c2-introduction]

\clearpage
\section{Notation Reference and Practice Procedure}
\label{sec:notation-and-procedure}
\phantomsection\label{tab:notation-reference}

\begingroup
\small
\setlength{\tabcolsep}{4pt}
\renewcommand{\arraystretch}{1.18}
\begin{longtable}{@{}>{\raggedright\arraybackslash}p{.17\linewidth}>{\raggedright\arraybackslash}p{.20\linewidth}>{\raggedright\arraybackslash}p{.28\linewidth}>{\raggedright\arraybackslash}p{\dimexpr.35\linewidth-24pt\relax}@{}}
\toprule
\textbf{Symbol or notation} & \textbf{Concept} & \textbf{Definition} & \textbf{Meaning in the introductory examples} \\
\midrule
\endfirsthead
\toprule
\textbf{Symbol or notation} & \textbf{Concept} & \textbf{Definition} & \textbf{Meaning in the introductory examples} \\
\midrule
\endhead
\midrule\multicolumn{4}{r}{\small Continued on the next page.}\\
\endfoot
\bottomrule
\endlastfoot
$\CtwoMean{\mu}$ & \CtwoMean{Population mean} & Fixed population parameter being estimated. & \textbullet\ \hyperref[subsec:c2-cycling-example]{Example 1}: unknown mean trip time, in minutes.\par \textbullet\ \hyperref[subsec:c2-coverage-example]{Example 2}: known simulation mean, \CtwoMean{41 minutes}. \\
$\CtwoPoint{\bar{x}}$ & \CtwoPoint{Point estimate} & Sample mean; a single estimate of $\CtwoMean{\mu}$. & \textbullet\ \hyperref[subsec:c2-cycling-example]{Example 1}: \CtwoPoint{40 minutes}.\par \textbullet\ \hyperref[subsec:c2-coverage-example]{Example 2}: recalculated for each sample. \\
$\CtwoSize{n}$ & \CtwoSize{Sample size} & Number of observations in a sample. & \textbullet\ \hyperref[subsec:c2-cycling-example]{Examples 1} and \hyperref[subsec:c2-coverage-example]{2}: \CtwoSize{49 trips} per sample. \\
$\CtwoSD{\sigma}$ & \CtwoSD{Population SD} & Known standard deviation of individual observations. & \textbullet\ \hyperref[subsec:c2-cycling-example]{Examples 1} and \hyperref[subsec:c2-coverage-example]{2}: \CtwoSD{14 minutes}. \\
$\CtwoSE{\mathrm{SE}}$ & \CtwoSE{Standard error} & Standard deviation of $\CtwoPoint{\bar{x}}$; $\CtwoSD{\sigma}/\sqrt{\CtwoSize{n}}$ for independent observations. & \textbullet\ \hyperref[subsec:c2-cycling-example]{Examples 1} and \hyperref[subsec:c2-coverage-example]{2}: \CtwoSE{2 minutes}, not the spread of individual trips. \\
$\CtwoConfidence{C}$ & \CtwoConfidence{Confidence level} & Chosen nominal long-run coverage, expressed as a decimal. & \textbullet\ \hyperref[subsec:c2-cycling-example]{Examples 1} and \hyperref[subsec:c2-coverage-example]{2}: \CtwoConfidence{0.95}, or \CtwoConfidence{95\%}. \\
$\CtwoSignificance{\alpha}$ & \CtwoSignificance{Significance level} & $1-\CtwoConfidence{C}$; total standard normal area outside the central region. & \textbullet\ \hyperref[subsec:c2-cycling-example]{Example 1}: \CtwoSignificance{0.05}.\par \textbullet\ \hyperref[subsec:c2-coverage-example]{Example 2}: unchanged by observed coverage. \\
$\CtwoSignificance{\alpha}/2$ & \CtwoSignificance{Tail area} & Half the significance level in each tail of a two-sided procedure. & \textbullet\ \hyperref[subsec:c2-cycling-example]{Example 1}: \CtwoSignificance{0.025} in each tail. \\
$\CtwoCritical{z^*}$ & \CtwoCritical{Critical value} & Unitless positive normal cutoff: $\Phi^{-1}(1-\CtwoSignificance{\alpha}/2)$. & \textbullet\ \hyperref[subsec:c2-cycling-example]{Example 1}: \CtwoCritical{1.960} for \CtwoConfidence{95\%}. \\
$\CtwoMoE{\mathrm{MoE}}$ & \CtwoMoE{Margin of error} & Distance from centre to either bound: $\CtwoCritical{z^*}\times\CtwoSE{\mathrm{SE}}$. & \textbullet\ \hyperref[subsec:c2-cycling-example]{Example 1}: \CtwoMoE{3.92 minutes}. \\
$\CtwoEndpoint{L}$, $\CtwoEndpoint{U}$;\par $\CtwoPoint{\bar{x}}\pm\CtwoMoE{\mathrm{MoE}}$ & \CtwoEndpoint{Endpoints} & Lower bound $\CtwoEndpoint{L}=\CtwoPoint{\bar{x}}-\CtwoMoE{\mathrm{MoE}}$; upper bound $\CtwoEndpoint{U}=\CtwoPoint{\bar{x}}+\CtwoMoE{\mathrm{MoE}}$. & \textbullet\ \hyperref[subsec:c2-cycling-example]{Example 1}: $[\CtwoEndpoint{36.08},\CtwoEndpoint{43.92}]$ minutes. \\
$[\CtwoEndpoint{L},\CtwoEndpoint{U}]$ & Interval estimate & Two-sided range of plausible values for $\CtwoMean{\mu}$ under the model. & \textbullet\ \hyperref[subsec:c2-cycling-example]{Example 1}: includes 42 minutes.\par \textbullet\ \hyperref[subsec:c2-coverage-example]{Example 2}: covers when $\CtwoEndpoint{L}\leq\CtwoMean{41}\leq\CtwoEndpoint{U}$. \\
$\CtwoEndpoint{U}-\CtwoEndpoint{L}$;\par $2\CtwoMoE{\mathrm{MoE}}$ & \CtwoWidth{Full width} & Upper minus lower bound; twice the margin of error. & \textbullet\ \hyperref[subsec:c2-cycling-example]{Example 1}: \CtwoWidth{7.84 minutes}. \\
$\CtwoCoverage{\text{rate}}$ & \CtwoCoverage{Observed coverage} & Covering count divided by total intervals for that procedure; multiply by $100\%$. & \textbullet\ \hyperref[subsec:c2-coverage-example]{Example 2}: 47 of 50, or \CtwoCoverage{94\%}; a finite-run result. \\
$\Phi$, $\Phi^{-1}$ & Normal area and inverse & $\Phi(z)$ gives area to the left of $z$; $\Phi^{-1}$ finds its cutoff. & \textbullet\ \hyperref[subsec:c2-comprehensive]{Worked example, Step 2}: left area $0.975$ selects $\CtwoCritical{1.960}$. \\
\end{longtable}
\endgroup

Greens connect sample quantities; reds connect population quantities. Light blue links confidence and significance, dark blue marks the critical value, and yellow marks standard error. Related orange shades distinguish margin of error, endpoints, and full width; purple marks observed coverage. Names and symbols remain the guide without color. See the \hyperref[tab:highlight-legend]{concise legend} for a quick reference.

\clearpage
\subsection{Practice procedure}

\begin{enumerate}
    \item \textbf{Task 1: identify, organize, calculate SE.} Make an information table with exactly the columns \textbf{Symbol}, \textbf{Concept}, \textbf{Type}, and \textbf{Meaning in context}, using one row per symbol and concept. Use the exact Type labels \textbf{Sample statistics} ($\bar{x}$, $s$ only when provided or relevant, $n$, $\mathrm{SE}$), \textbf{Population parameters} ($\mu$, $\sigma$), and \textbf{Procedure Settings} ($\CtwoConfidence{C}$, $\CtwoSignificance{\alpha}$, $\CtwoSignificance{\alpha}/2$, $\CtwoCritical{z^*}$). Classify $n$ and $\mathrm{SE}$ under Sample statistics by this organizational convention. Keep notation, concept names, and classifications in their respective columns; put contextual values, units, interpretations, and calculations in Meaning in context. Explain the unknown population mean and its point estimate. Label values by team or confidence level when comparing procedures. Record $\CtwoConfidence{C}$ as a decimal and give $\CtwoSignificance{\alpha}=1-\CtwoConfidence{C}$, $\CtwoSignificance{\alpha}/2$, and $\CtwoCritical{z^*}=\Phi^{-1}(1-\CtwoSignificance{\alpha}/2)$ in separate rows. Calculate $\CtwoSE{\mathrm{SE}}=\CtwoSD{\sigma}/\sqrt{\CtwoSize{n}}$ with substitution, result, and units for each sample size in the SE row.
    \item \textbf{Task 2: calculate MoE and state the interval.} Use $\CtwoMoE{\mathrm{MoE}}=\CtwoCritical{z^*}\CtwoSE{\mathrm{SE}}$, then $\CtwoPoint{\bar{x}}\pm\CtwoMoE{\mathrm{MoE}}$. Interpret the range as plausible values for the population mean under the model. Compare widths $2\CtwoMoE{\mathrm{MoE}}$. If endpoints are rounded, distinguish the unrounded width from the reported width (reported upper minus reported lower endpoint).
    \item \textbf{Task 3: calculate observed coverage and choose one option.} Extract the numerator and denominator from the narrative, diagram, or table; use the coverage formula below separately for each procedure. Explain the choice using the distinction between finite-run observed coverage and nominal long-run confidence.
\end{enumerate}


Use the same units for $\CtwoMean{\mu}$, $\CtwoPoint{\bar{x}}$, $\CtwoSD{\sigma}$, $\CtwoSE{\mathrm{SE}}$, $\CtwoMoE{\mathrm{MoE}}$, the \CtwoEndpoint{endpoints}, and the \CtwoWidth{width}. Population standard deviations shown with $\approx$ are rounded to two decimal places; calculations retain their unrounded values.

\subsection{Observed coverage and the fixed mean}
\[
\CtwoCoverage{\text{Observed coverage}}=
\frac{\text{number of intervals containing the fixed population mean}}
{\text{total number of simulated intervals}}\times100\%.
\]
In the original estimation problem, $\CtwoMean{\mu}$ is fixed but unknown. In a simulation, the model supplies a known fixed $\CtwoMean{\mu}$ so that coverage can be checked: an interval $[L,U]$ covers exactly when $L\leq\mu\leq U$, including equality at either endpoint.

For \textbf{narrative evidence}, use the covering count, or subtract misses from the total. For a \textbf{diagram}, count the horizontal intervals meeting or crossing the fixed-mean reference line, then count all intervals. For a \textbf{table}, sum the relevant covering cells and use the total number of intervals for that procedure. Paired procedures applied to 100 samples each have denominator 100, even though there are 200 intervals altogether.

The \CtwoConfidence{nominal confidence level} $\CtwoConfidence{C}$ is chosen before sampling and sets $\CtwoSignificance{\alpha}=1-\CtwoConfidence{C}$. Never calculate $\CtwoSignificance{\alpha}$ from \CtwoCoverage{observed coverage}. Under the model, $\CtwoSignificance{\alpha}$ is the long-run fraction of intervals that miss the fixed mean, not a probability that a particular realized interval misses it. \CtwoCoverage{Observed coverage} is calculated after a finite run and can differ from that level. Under the stated assumptions, confidence describes the procedure's long-run coverage of the fixed parameter. A realized interval either contains the parameter or does not; 95\% confidence does not assign a 95\% probability to that fixed interval containing it, and does not describe the percentage of individual observations inside it.

\clearpage
\subsection{Choosing the critical value}

Write the \CtwoConfidence{nominal confidence level} $\CtwoConfidence{C}$ as a decimal and define the \CtwoSignificance{significance level} by $\CtwoSignificance{\alpha}=1-\CtwoConfidence{C}$. For a two-sided confidence interval, the standard normal curve has \CtwoConfidence{central area} $\CtwoConfidence{C}$ between $-\CtwoCritical{z^*}$ and $\CtwoCritical{z^*}$, leaving area $\CtwoSignificance{\alpha}/2$ in each tail. The area to the left of $\CtwoCritical{z^*}$ is $1-\CtwoSignificance{\alpha}/2$, so
\[
    \CtwoCritical{z^*}=\Phi^{-1}\!\left(1-\CtwoSignificance{\alpha}/2\right).
\]
\noindent
\begin{minipage}[t]{0.56\linewidth}
    \vspace{0pt}
    $\Phi$ gives the area to the left of a value on the standard normal curve. Its inverse, $\Phi^{-1}$, takes an area and finds the value with that much area to its left. For \CtwoConfidence{95\%} confidence, $\CtwoSignificance{\alpha}=\CtwoSignificance{0.05}$ and each tail has area $\CtwoSignificance{\alpha}/2=\CtwoSignificance{0.025}$, so look for $1-\CtwoSignificance{0.025}=0.975$ to the left of $\CtwoCritical{z^*}$.

    \smallskip
    Choose $\CtwoCritical{z^*}$ from the \CtwoConfidence{confidence level} and calculate $\CtwoPoint{\bar{x}}$ from the sample data. \CtwoConfidence{Higher confidence} means smaller $\CtwoSignificance{\alpha}$ and requires a larger $\CtwoCritical{z^*}$ to enclose more area. With the \CtwoSE{same standard error}, this gives a larger $\CtwoMoE{\mathrm{MoE}}$ and a \CtwoWidth{wider interval}.
\end{minipage}
\hfill
\begin{minipage}[t]{0.40\linewidth}
    \vspace{0pt}
    \centering
    \begin{tabular}{@{}cccc@{}}
        \toprule
        $\CtwoConfidence{C}$ & $\CtwoSignificance{\alpha}$ & $\CtwoSignificance{\alpha}/2$ & $\CtwoCritical{z^*}$ \\
        \midrule
        \CtwoConfidence{80\%} & \CtwoSignificance{0.20} & \CtwoSignificance{0.10} & \CtwoCritical{1.282} \\
        \CtwoConfidence{90\%} & \CtwoSignificance{0.10} & \CtwoSignificance{0.05} & \CtwoCritical{1.645} \\
        \CtwoConfidence{95\%} & \CtwoSignificance{0.05} & \CtwoSignificance{0.025} & \CtwoCritical{1.960} \\
        \CtwoConfidence{97.5\%} & \CtwoSignificance{0.025} & \CtwoSignificance{0.0125} & \CtwoCritical{2.241} \\
        \CtwoConfidence{99\%} & \CtwoSignificance{0.01} & \CtwoSignificance{0.005} & \CtwoCritical{2.576} \\
        \CtwoConfidence{99.9\%} & \CtwoSignificance{0.001} & \CtwoSignificance{0.0005} & \CtwoCritical{3.291} \\
        \bottomrule
    \end{tabular}
    \par\smallskip
    \small Two-sided values rounded to three decimals.
\end{minipage}

\smallskip
\noindent
\begin{minipage}{\linewidth}
    \centering
    \begin{tikzpicture}
    \begin{axis}[
        width=0.82\linewidth,height=4.8cm,axis x line=bottom,axis y line=none,
        xmin=-4.2,xmax=4.2,ymin=0,ymax=0.43,
        xlabel={$z$},xlabel style={at={(axis description cs:1,0)},anchor=west,font=\small},
        xtick={-1.96,0,1.96},xticklabels={$-\CtwoCritical{z^*}=-\CtwoCritical{1.960}$,$0$,$\CtwoCritical{z^*}=\CtwoCritical{1.960}$},
        tick label style={font=\small},samples=101,
        declare function={normaldensity(\x)=exp(-\x^2/2)/sqrt(2*pi);},
        ]
        \addplot[draw=none,fill=black!12,domain=-4:-1.96] {normaldensity(x)} \closedcycle;
        \addplot[draw=none,fill=CtwoCentralArea,domain=-1.96:1.96] {normaldensity(x)} \closedcycle;
        \addplot[draw=none,fill=black!12,domain=1.96:4] {normaldensity(x)} \closedcycle;
        \addplot[black,thick,smooth,domain=-4:4] {normaldensity(x)};
        \addplot[black,densely dashed] coordinates {(-1.96,0) (-1.96,{normaldensity(1.96)})};
        \addplot[black,densely dashed] coordinates {(1.96,0) (1.96,{normaldensity(1.96)})};
        \node[font=\small] at (axis cs:0,0.17) {$\CtwoConfidence{C}=\CtwoConfidence{0.95}$};
        \node[font=\small] (left tail) at (axis cs:-2.8,0.20) {$\CtwoSignificance{\alpha}/2=\CtwoSignificance{0.025}$};
        \node[font=\small] (right tail) at (axis cs:2.8,0.20) {$\CtwoSignificance{\alpha}/2=\CtwoSignificance{0.025}$};
        \draw[->] (left tail.south) -- (axis cs:-2.18,0.017);
        \draw[->] (right tail.south) -- (axis cs:2.18,0.017);
    \end{axis}
    \end{tikzpicture}
    \par\smallskip
    \small The cumulative area to the left of the axis position $\CtwoCritical{z^*}=\CtwoCritical{1.960}$ is the left tail plus the central region, so $\Phi(\CtwoCritical{1.960})\approx \CtwoSignificance{0.025}+\CtwoConfidence{0.95}=1-\CtwoSignificance{\alpha}/2=0.975$. Here $\CtwoSignificance{\alpha}=\CtwoSignificance{0.05}$ is split equally between the two tails.
\end{minipage}

\clearpage
\section{General Worked Problems}
\label{sec:general-worked}

\SectionProblemsTableOfContents{
    \SectionProblemLink{prob:c2-1}{Repeated estimates of market fees} \\
    \SectionProblemLink{prob:c2-2}{Two reports on household energy bills} \\
    \SectionProblemLink{prob:c2-3}{Bus commute times in a simulation} \\
    \SectionProblemLink{prob:c2-compare-three}{Three confidence levels for repair times}
}

\subsection{Repeated estimates of market fees}
\label{prob:c2-1}

\noindent
\begin{minipage}[t]{0.48\linewidth}
    \vspace{0pt}
    A city takes a random sample of $n=100$ market stalls to estimate the mean weekly stall fee. In thousand COP, the sample mean is $\bar{x}=106$ and the known population standard deviation is $\sigma\approx25.51$. Assume independent observations and normally distributed sample means. The nominal confidence level $C=0.95$ (95\%) sets $\alpha=1-C$ and uses the standard normal critical value $z^*=1.960$.

    The diagram shows the 95\% intervals from five fresh samples, each of size $n=100$, from a simulated market with known fixed mean $\mu=104$ thousand COP and the same population standard deviation. The dashed line marks that mean; each horizontal segment is one interval.
\end{minipage}
\hfill
\begin{minipage}[t]{0.48\linewidth}
    \vspace{0pt}
    \centering
    \begin{tikzpicture}
    \begin{axis}[
        width=\linewidth,height=4.6cm,axis lines=left,
        xmin=98,xmax=116,ymin=0.6,ymax=5.9,
        xlabel={Mean weekly fee (thousand COP)},ylabel={Simulated sample},
        ytick={1,2,3,4,5},xtick={99,104,109,114},
        tick label style={font=\small},label style={font=\small}]
        \addplot[black,dashed] coordinates {(104,0.6) (104,5.65)};
        \addplot[black,thick,mark=|] coordinates {(101,1) (111,1)};
        \addplot[black,thick,mark=|] coordinates {(99,2) (109,2)};
        \addplot[black,thick,mark=|] coordinates {(103,3) (113,3)};
        \addplot[black,thick,mark=|] coordinates {(105,4) (115,4)};
        \addplot[black,thick,mark=|] coordinates {(102,5) (112,5)};
        \node[font=\small,anchor=east] at (axis cs:103.7,4.4) {$\mu=104$};
    \end{axis}
    \end{tikzpicture}
\end{minipage}

\QuestionsTasks[sec:general-worked]{
    {2}/{sol:c2-1-a}/{\InformationTask},
    {2}/{sol:c2-1-b}/{\IntervalTask},
    {2}/{sol:c2-1-c}/{\CoverageTask\ Using the diagram, which option gives the correct number of intervals that cover the simulated mean and the correct interpretation of the city's 95\% confidence level?}
}

\begin{answeroptions}
    \item Four of five cover, so the observed 80\% is the city's long-run coverage rate.
    \item Four of five cover. The procedure covers the mean about 95\% of the time in the long run under the model. The city's calculated interval either covers its fixed mean or misses it.
    \item Four of five cover. The city's calculated interval has a 95\% probability of covering its fixed mean.
    \item Five of five cover. The procedure covers the mean about 95\% of the time in the long run under the model. The city's calculated interval either covers its fixed mean or misses it.
\end{answeroptions}

\clearpage
\subsubsection*{Task 1 solution. Organize the information and calculate SE}
\label{sol:c2-1-a}

\InformationTable{
    $\CtwoMean{\mu}$ & \CtwoMean{Population mean} & Population parameters & The mean weekly stall fee for all market stalls in the city, in \CtwoMean{thousand COP}. The fixed, unknown mean the city wants to estimate. \\
    $\CtwoPoint{\bar{x}}$ & \CtwoPoint{Sample mean} & Sample statistics & $\CtwoPoint{\bar{x}}=\CtwoPoint{106\,\text{thousand COP}}$. A sample statistic giving a single estimate of the city's mean weekly stall fee. \\
    $\CtwoSize{n}$ & \CtwoSize{Sample size} & Sample statistics & $\CtwoSize{n}=\CtwoSize{100\,\text{stalls}}$. Number sampled; not a population parameter. \\
    $\CtwoSD{\sigma}$ & \CtwoSD{Population} \CtwoSD{standard} \CtwoSD{deviation} & Population parameters & $\CtwoSD{\sigma}\approx\CtwoSD{25.51\,\text{thousand COP}}$. Known population parameter. \\
    $\CtwoConfidence{C}$ & \CtwoConfidence{Confidence level} & Procedure Settings & \CtwoConfidence{95\%} procedure: $\CtwoConfidence{C}=\CtwoConfidence{0.95}$. Nominal confidence level. \\
    $\CtwoSignificance{\alpha}$ & \CtwoSignificance{Significance level} & Procedure Settings & \CtwoConfidence{95\%} procedure: $\CtwoSignificance{\alpha}=1-\CtwoConfidence{0.95}=\CtwoSignificance{0.05}$. Total area outside the central region. \\
    $\CtwoSignificance{\alpha}/2$ & \CtwoSignificance{Tail area} & Procedure Settings & \CtwoConfidence{95\%} procedure: $\CtwoSignificance{\alpha}/2=\CtwoSignificance{0.025}$ in each tail. Area in each tail of the two-sided procedure. \\
    $\CtwoCritical{z^*}$ & \CtwoCritical{Critical value} & Procedure Settings & \CtwoConfidence{95\%} procedure: $\CtwoCritical{z^*}=\CtwoCritical{1.960}$. Corresponding standard normal critical value. \\
    $\CtwoSE{\mathrm{SE}}$ & \CtwoSE{Standard error} & Sample statistics & $\CtwoSE{\mathrm{SE}}=\CtwoSD{\sigma}/\sqrt{\CtwoSize{n}}\approx\CtwoSD{25.51}/\sqrt{\CtwoSize{100}}\approx\CtwoSE{125/49\,\text{thousand COP}}$. The standard deviation of the \CtwoPoint{sample mean} across repeated samples. \\
}

\par\vspace{3mm}\secInicio[prob:c2-1][sec:general-worked]

\subsubsection*{Task 2 solution. Calculate MoE and state the interval}
\label{sol:c2-1-b}

For each procedure, $\CtwoMoE{\mathrm{MoE}}=\CtwoCritical{z^*}\times\CtwoSE{\mathrm{SE}}$ and the interval is $[\CtwoPoint{\bar{x}}-\CtwoMoE{\mathrm{MoE}},\CtwoPoint{\bar{x}}+\CtwoMoE{\mathrm{MoE}}]$.

\EstimationTable{
    \CtwoConfidence{95\%} interval & $\CtwoMoE{\mathrm{MoE}}=\CtwoCritical{z^*}\CtwoSE{\mathrm{SE}}=\CtwoCritical{1.960}(\CtwoSE{125/49})=\CtwoMoE{5\,\text{thousand COP}}$. The interval is $[\CtwoPoint{106}-\CtwoMoE{5},\CtwoPoint{106}+\CtwoMoE{5}]=[\CtwoEndpoint{101},\CtwoEndpoint{111}]$ \CtwoEndpoint{thousand COP}. & A range of plausible values for the city's mean weekly stall fee under the model. \\
}

\par\vspace{3mm}\secInicio[prob:c2-1][sec:general-worked]


\noindent\begin{minipage}{\linewidth}
\subsubsection*{Task 3 solution. Calculate coverage and interpret confidence}
\label{sol:c2-1-c}

\textbf{Correct option B.} Each simulated interval is \CtwoWidth{10 thousand COP} wide. Samples 1, 2, 3, and 5 cross the line at \CtwoMean{104}. Sample 4 starts at \CtwoEndpoint{105} and misses, so the \CtwoCoverage{observed coverage} is $(4/5)\times100\%=\CtwoCoverage{80\%}$. The \CtwoConfidence{nominal level} is \CtwoConfidence{95\%}. Over many fresh samples under the model assumptions, about \CtwoConfidence{95\%} of the resulting intervals would contain the fixed mean.

The claim of a 95\% probability for the city's interval fails. Its \CtwoEndpoint{endpoints}, \CtwoEndpoint{101} and \CtwoEndpoint{111}, are fixed after sampling. The city's unknown mean either \CtwoCoverage{lies between them} or \CtwoCoverage{falls outside them}.

\par\vspace{3mm}\secInicio[prob:c2-1][sec:general-worked]
\end{minipage}\par

\clearpage

\subsection{Two reports on household energy bills}
\label{prob:c2-2}

A utility takes a random sample of $n=100$ household accounts from one billing month. Both reports use this sample, with mean $\bar{x}=120$ thousand COP. The known population standard deviation is $\sigma=38$ thousand COP. Assume independent observations and a normal sampling model. The table gives the nominal confidence level $C$ and standard normal critical value for each report. Identify $\alpha=1-C$ and each tail area $\alpha/2$ in Task 1.

\begin{center}
\begin{tabular}{@{}lcc@{}}
    \toprule
    Nominal confidence level & Critical value $z^*$ & Interval width (thousand COP) \\
    \midrule
    90\% & 1.645 & \\
    99\% & 2.576 & \\
    \bottomrule
\end{tabular}
\end{center}

The reports round endpoints to the nearest thousand COP. The confidence levels refer to the procedures before rounding.

In a separate, explicitly hypothetical simulation, a normal model has known fixed mean $\mu=120$ thousand COP and the same $\sigma=38$. From each of 100 fresh samples of size $n=100$, the utility constructs both an unrounded 90\% interval and an unrounded 99\% interval using the stated critical values. Of the 100 unrounded 90\% intervals, 11 miss the fixed mean; of the 100 unrounded 99\% intervals, 2 miss it. Coverage is evaluated before rounding, separately for each level. These simulated samples are not the original household sample.


\QuestionsTasks[sec:general-worked]{
    {2}/{sol:c2-2-a}/{\textbf{Task 1 -- Organize the information and calculate SE.}},
    {2}/{sol:c2-2-b}/{\textbf{Task 2 -- Calculate MoE and state the interval.} Use endpoints rounded to the nearest thousand COP to complete the width table; also state the unrounded endpoints and widths.},
    {2}/{sol:c2-2-c}/{\textbf{Task 3 --} Which option gives observed coverage for the 90\% and 99\% procedures, in that order, correctly compares the widths, and interprets 99\% confidence?}
}

\begin{answeroptions}
    \item The observed rates are 89\% and 98\%, respectively; the 99\% interval is wider and contains 99\% of individual bills.
    \item The observed rates are 90\% and 99\%, respectively; each finite run must match its nominal confidence level.
    \item The observed rates are 89\% and 98\%, respectively; the wider reported interval has a 99\% probability of containing the fixed mean.
    \item The observed rates are 89\% and 98\%, respectively. The 99\% interval is wider; its unrounded procedure covers the fixed mean about 99\% of the time in the long run under the model, without specifying a percentage of individual bills inside it.
\end{answeroptions}

\clearpage
\subsubsection*{Task 1 solution. Organize the information and calculate SE}
\label{sol:c2-2-a}

\InformationTable{
    $\CtwoMean{\mu}$ & \CtwoMean{Population mean} & Population parameters & The mean energy bill for all the utility's household accounts in that billing month, in \CtwoMean{thousand COP}. The fixed, unknown mean estimated in both reports. \\
    $\CtwoPoint{\bar{x}}$ & \CtwoPoint{Sample mean} & Sample statistics & $\CtwoPoint{\bar{x}}=\CtwoPoint{120\,\text{thousand COP}}$. A sample statistic giving a single estimate of the mean household bill for the month. \\
    $\CtwoSize{n}$ & \CtwoSize{Sample size} & Sample statistics & $\CtwoSize{n}=\CtwoSize{100\,\text{accounts}}$, shared. Both reports use the same sample. \\
    $\CtwoSD{\sigma}$ & \CtwoSD{Population} \CtwoSD{standard} \CtwoSD{deviation} & Population parameters & $\CtwoSD{\sigma}=\CtwoSD{38\,\text{thousand COP}}$, shared. Known population parameter. \\
    $\CtwoConfidence{C}$ & \CtwoConfidence{Confidence level} & Procedure Settings & \CtwoConfidence{90\%} procedure: $\CtwoConfidence{C}=\CtwoConfidence{0.90}$.\par \CtwoConfidence{99\%} procedure: $\CtwoConfidence{C}=\CtwoConfidence{0.99}$. Nominal confidence level. \\
    $\CtwoSignificance{\alpha}$ & \CtwoSignificance{Significance level} & Procedure Settings & \CtwoConfidence{90\%} procedure: $\CtwoSignificance{\alpha}=1-\CtwoConfidence{0.90}=\CtwoSignificance{0.10}$.\par \CtwoConfidence{99\%} procedure: $\CtwoSignificance{\alpha}=1-\CtwoConfidence{0.99}=\CtwoSignificance{0.01}$. Total area outside the central region. \\
    $\CtwoSignificance{\alpha}/2$ & \CtwoSignificance{Tail area} & Procedure Settings & \CtwoConfidence{90\%} procedure: $\CtwoSignificance{\alpha}/2=\CtwoSignificance{0.05}$ in each tail.\par \CtwoConfidence{99\%} procedure: $\CtwoSignificance{\alpha}/2=\CtwoSignificance{0.005}$ in each tail. Area in each tail of the two-sided procedure. \\
    $\CtwoCritical{z^*}$ & \CtwoCritical{Critical value} & Procedure Settings & \CtwoConfidence{90\%} procedure: $\CtwoCritical{z^*}=\CtwoCritical{1.645}$.\par \CtwoConfidence{99\%} procedure: $\CtwoCritical{z^*}=\CtwoCritical{2.576}$. Corresponding standard normal critical value. \\
    $\CtwoSE{\mathrm{SE}}$ & \CtwoSE{Standard error} & Sample statistics & $\CtwoSE{\mathrm{SE}}=\CtwoSD{\sigma}/\sqrt{\CtwoSize{n}}=\CtwoSD{38}/\sqrt{\CtwoSize{100}}=\CtwoSE{3.8\,\text{thousand COP}}$. Both reports use the same sample and \CtwoSE{standard error}. \\
}

\par\vspace{3mm}\secInicio[prob:c2-2][sec:general-worked]

\clearpage
\subsubsection*{Task 2 solution. Calculate MoE and state the interval}
\label{sol:c2-2-b}

For each procedure, $\CtwoMoE{\mathrm{MoE}}=\CtwoCritical{z^*}\times\CtwoSE{\mathrm{SE}}$ and the interval is $[\CtwoPoint{\bar{x}}-\CtwoMoE{\mathrm{MoE}},\CtwoPoint{\bar{x}}+\CtwoMoE{\mathrm{MoE}}]$.

\EstimationTable{
    \CtwoConfidence{90\%} interval & $\CtwoMoE{\mathrm{MoE}}=\CtwoCritical{z^*}\CtwoSE{\mathrm{SE}}=\CtwoCritical{1.645}(\CtwoSE{3.8})=\CtwoMoE{6.251\,\text{thousand COP}}$. The \CtwoEndpoint{endpoints} from $\CtwoPoint{120}\pm\CtwoMoE{6.251}$ are $\CtwoEndpoint{113.749}$ and $\CtwoEndpoint{126.251}$. The rounded interval is $[\CtwoEndpoint{114},\CtwoEndpoint{126}]$, with \CtwoWidth{width} $\CtwoEndpoint{126}-\CtwoEndpoint{114}=\CtwoWidth{12\,\text{thousand COP}}$. & The narrower reported range of plausible values for the mean household bill under the model. \\
    \CtwoConfidence{99\%} interval & $\CtwoMoE{\mathrm{MoE}}=\CtwoCritical{z^*}\CtwoSE{\mathrm{SE}}=\CtwoCritical{2.576}(\CtwoSE{3.8})=\CtwoMoE{9.7888\,\text{thousand COP}}$. The \CtwoEndpoint{endpoints} from $\CtwoPoint{120}\pm\CtwoMoE{9.7888}$ are $\CtwoEndpoint{110.2112}$ and $\CtwoEndpoint{129.7888}$. The rounded interval is $[\CtwoEndpoint{110},\CtwoEndpoint{130}]$, with \CtwoWidth{width} $\CtwoEndpoint{130}-\CtwoEndpoint{110}=\CtwoWidth{20\,\text{thousand COP}}$. & The wider reported range of plausible values for the same mean under the model. \\
}

Before rounding, the \CtwoWidth{widths} are $2(\CtwoMoE{6.251})=\CtwoWidth{12.502}$ and $2(\CtwoMoE{9.7888})=\CtwoWidth{19.5776\,\text{thousand COP}}$. The reported \CtwoWidth{widths}, in \CtwoWidth{thousand COP}, are $\CtwoEndpoint{126}-\CtwoEndpoint{114}=\CtwoWidth{12}$ and $\CtwoEndpoint{130}-\CtwoEndpoint{110}=\CtwoWidth{20}$. With the \CtwoSE{same standard error}, the \CtwoCritical{larger critical value} for \CtwoConfidence{99\%} gives a larger $\CtwoMoE{\mathrm{MoE}}$ and puts the \CtwoEndpoint{endpoints} farther from \CtwoPoint{120}.

\begin{center}
\begin{tikzpicture}
\begin{axis}[
    width=0.88\linewidth,height=4.3cm,axis lines=left,
    xmin=108,xmax=132,ymin=0.35,ymax=2.7,
    xlabel={Mean monthly household energy bill (thousand COP)},
    ytick={1,2},
    yticklabels={\CtwoConfidence{90\%},\CtwoConfidence{99\%}},
    xtick={110,115,120,125,130},tick label style={font=\small},label style={font=\small}]
    \addplot[black,dashed] coordinates {(120,0.35) (120,2.3)};
    \node[anchor=south,font=\small] at (axis cs:120,2.3) {$\CtwoPoint{\bar{x}}$};
    \addplot[black,thick,mark=|,mark size=5pt] coordinates {(113.749,1) (126.251,1)};
    \addplot[black,only marks,mark=*,mark size=2pt] coordinates {(120,1)};
    \node[anchor=south,font=\small] at (axis cs:113.749,1.05) {$\CtwoEndpoint{113.749}$};
    \node[anchor=south,font=\small] at (axis cs:126.251,1.05) {$\CtwoEndpoint{126.251}$};
    \addplot[black,thick,mark=|,mark size=5pt] coordinates {(110.2112,2) (129.7888,2)};
    \addplot[black,only marks,mark=*,mark size=2pt] coordinates {(120,2)};
    \node[anchor=south,font=\small] at (axis cs:110.2112,2.05) {$\CtwoEndpoint{110.2112}$};
    \node[anchor=south,font=\small] at (axis cs:129.7888,2.05) {$\CtwoEndpoint{129.7888}$};
\end{axis}
\end{tikzpicture}
\end{center}

This graph compares the two unrounded intervals from the same sample. Both centre markers and the dashed line mark the \CtwoPoint{sample mean} $\CtwoPoint{\bar{x}}=\CtwoPoint{120\,\text{thousand COP}}$; the actual \CtwoMean{population mean} $\CtwoMean{\mu}$ remains unknown. Increasing confidence widens the interval while keeping this centre fixed. The graph uses the calculated endpoints before rounding. The reports instead round each endpoint to the nearest thousand COP and calculate the reported width from those rounded endpoints; nominal confidence and simulated coverage refer to the unrounded procedures.

\par\vspace{3mm}\secInicio[prob:c2-2][sec:general-worked]


\noindent\begin{minipage}{\linewidth}
\subsubsection*{Task 3 solution. Calculate coverage and interpret confidence}
\label{sol:c2-2-c}

\textbf{Correct option D.} There are $100-11=89$ covering 90\% intervals out of 100, so $(89/100)\times100\%=\CtwoCoverage{89\%}$. There are $100-2=98$ covering 99\% intervals out of 100, so $(98/100)\times100\%=\CtwoCoverage{98\%}$. Each denominator is 100, not 200; both counts concern unrounded intervals. The finite-run rates need not equal the nominal levels.

Under the model assumptions, about \CtwoConfidence{90\%} and \CtwoConfidence{99\%} of the respective unrounded intervals would cover the fixed mean over repeated samples. Each calculated interval either \CtwoCoverage{covers the mean} or \CtwoCoverage{misses it}. The claim about 99\% of individual bills fails because coverage concerns the \CtwoMean{population mean}. The proportion of bills from \CtwoEndpoint{110} to \CtwoEndpoint{130 thousand COP} depends on the distribution of individual bills.

\par\vspace{3mm}\secInicio[prob:c2-2][sec:general-worked]
\end{minipage}\par

\clearpage

\subsection{Bus commute times in a simulation}
\label{prob:c2-3}

A transit agency takes an independent random sample of $n=36$ weekday passenger trips on one bus route. The sample mean is $\bar{x}=34$ minutes. Assume normally distributed trip times with known population standard deviation $\sigma=12$ minutes. The nominal confidence level $C=0.90$ (90\%) sets $\alpha=1-C$ and uses the standard normal critical value $z^*=1.645$.

To study the 90\% procedure, the agency simulates 20 fresh samples, each of size $n=36$, from a normal model with mean 35 minutes and the same standard deviation. The table groups the simulated intervals by sample number.

\begin{center}
\begin{tabular}{@{}lcccc@{}}
    \toprule
    Sample numbers & 1--5 & 6--10 & 11--15 & 16--20 \\
    \midrule
    Intervals containing 35 minutes & 5 & 4 & 3 & 5 \\
    \bottomrule
\end{tabular}
\end{center}

\QuestionsTasks[sec:general-worked]{
    {2}/{sol:c2-3-a}/{\textbf{Task 1 -- Organize the information and calculate SE.}},
    {2}/{sol:c2-3-b}/{\textbf{Task 2 -- Calculate MoE and state the interval.}},
    {2}/{sol:c2-3-c}/{\textbf{Task 3 --} Which option correctly totals the covering intervals in the table and assesses the claim that every run of 20 must have exactly 18 intervals that cover the mean?}
}

\begin{answeroptions}
    \item Seventeen of 20 cover. The claim fails because the 90\% procedure can give different counts in different runs of 20.
    \item Eighteen of 20 cover. The claim fails because the 90\% procedure can give different counts in different runs of 20.
    \item Seventeen of 20 cover. The claim fails because the observed 85\% is the procedure's long-run coverage rate.
    \item Seventeen of 20 cover. The claim fails because each calculated interval has a 90\% probability of containing the fixed mean.
\end{answeroptions}

\clearpage
\subsubsection*{Task 1 solution. Organize the information and calculate SE}
\label{sol:c2-3-a}

\InformationTable{
    $\CtwoMean{\mu}$ & \CtwoMean{Population mean} & Population parameters & The mean duration of all weekday passenger trips on the route, in \CtwoMean{minutes}. The fixed, unknown mean the transit agency wants to estimate. \\
    $\CtwoPoint{\bar{x}}$ & \CtwoPoint{Sample mean} & Sample statistics & $\CtwoPoint{\bar{x}}=\CtwoPoint{34\,\text{minutes}}$. A sample statistic giving a single estimate of the mean weekday trip time on the route. \\
    $\CtwoSize{n}$ & \CtwoSize{Sample size} & Sample statistics & $\CtwoSize{n}=\CtwoSize{36\,\text{passenger trips}}$. Number sampled. \\
    $\CtwoSD{\sigma}$ & \CtwoSD{Population} \CtwoSD{standard} \CtwoSD{deviation} & Population parameters & $\CtwoSD{\sigma}=\CtwoSD{12\,\text{minutes}}$. Known population parameter. \\
    $\CtwoConfidence{C}$ & \CtwoConfidence{Confidence level} & Procedure Settings & \CtwoConfidence{90\%} procedure: $\CtwoConfidence{C}=\CtwoConfidence{0.90}$. Nominal confidence level. \\
    $\CtwoSignificance{\alpha}$ & \CtwoSignificance{Significance level} & Procedure Settings & \CtwoConfidence{90\%} procedure: $\CtwoSignificance{\alpha}=1-\CtwoConfidence{0.90}=\CtwoSignificance{0.10}$. Total area outside the central region. \\
    $\CtwoSignificance{\alpha}/2$ & \CtwoSignificance{Tail area} & Procedure Settings & \CtwoConfidence{90\%} procedure: $\CtwoSignificance{\alpha}/2=\CtwoSignificance{0.05}$ in each tail. Area in each tail of the two-sided procedure. \\
    $\CtwoCritical{z^*}$ & \CtwoCritical{Critical value} & Procedure Settings & \CtwoConfidence{90\%} procedure: $\CtwoCritical{z^*}=\CtwoCritical{1.645}$. Corresponding standard normal critical value. \\
    $\CtwoSE{\mathrm{SE}}$ & \CtwoSE{Standard error} & Sample statistics & $\CtwoSE{\mathrm{SE}}=\CtwoSD{\sigma}/\sqrt{\CtwoSize{n}}=\CtwoSD{12}/\sqrt{\CtwoSize{36}}=\CtwoSE{2\,\text{minutes}}$. The standard deviation of the \CtwoPoint{sample mean} across repeated samples. \\
}

\par\vspace{3mm}\secInicio[prob:c2-3][sec:general-worked]

\subsubsection*{Task 2 solution. Calculate MoE and state the interval}
\label{sol:c2-3-b}

For each procedure, $\CtwoMoE{\mathrm{MoE}}=\CtwoCritical{z^*}\times\CtwoSE{\mathrm{SE}}$ and the interval is $[\CtwoPoint{\bar{x}}-\CtwoMoE{\mathrm{MoE}},\CtwoPoint{\bar{x}}+\CtwoMoE{\mathrm{MoE}}]$.

\EstimationTable{
    \CtwoConfidence{90\%} interval & $\CtwoMoE{\mathrm{MoE}}=\CtwoCritical{z^*}\CtwoSE{\mathrm{SE}}=\CtwoCritical{1.645}(\CtwoSE{2})=\CtwoMoE{3.29\,\text{minutes}}$. The \CtwoEndpoint{endpoints} from $\CtwoPoint{34}\pm\CtwoMoE{3.29}$ give $[\CtwoEndpoint{30.71},\CtwoEndpoint{37.29}]$ \CtwoEndpoint{minutes}. & A range of plausible values for the route's mean weekday trip time under the model. \\
}

\par\vspace{3mm}\secInicio[prob:c2-3][sec:general-worked]


\noindent\begin{minipage}{\linewidth}
\subsubsection*{Task 3 solution. Calculate coverage and interpret confidence}
\label{sol:c2-3-c}

\textbf{Correct option A.} The table gives $5+4+3+5=17$ intervals that contain \CtwoMean{35 minutes}, for \CtwoCoverage{observed coverage} of $(17/20)\times100\%=\CtwoCoverage{85\%}$. This run disproves the claim of exactly 18 covering intervals in every run.

The \CtwoConfidence{nominal level} is \CtwoConfidence{90\%}. Under the model assumptions, about \CtwoConfidence{90\%} of intervals cover the fixed mean over many repetitions. This gives an average of about $\CtwoConfidence{0.90}(20)=18$ covering intervals per run of 20, with counts that can vary from run to run. Each calculated interval either \CtwoCoverage{contains the mean} or \CtwoCoverage{misses it}.

\par\vspace{3mm}\secInicio[prob:c2-3][sec:general-worked]
\end{minipage}\par

\clearpage

\subsection{Three confidence levels for repair times}
\label{prob:c2-compare-three}

A repair cooperative takes one independent random sample of $n=64$ completed appliance repairs from a specified month. The sample mean repair time is $\bar{x}=50$ minutes. Assume normally distributed repair times with known population standard deviation $\sigma=16$ minutes. Use this same sample to calculate exactly three two-sided confidence intervals: 90\%, 95\%, and 99\%, with $z^*=1.645$, $1.960$, and $2.576$, respectively. Keep the sample mean, sample size, population standard deviation, and standard error unchanged across the three procedures. Do not round intermediate calculations or endpoints.

In a separate hypothetical simulation, a normal population has known fixed mean $\mu=51$ minutes and the same $\sigma=16$ minutes. From each of 200 fresh samples of size $n=64$, all three unrounded intervals are constructed around that sample's mean. The mutually exclusive rows below count the outcomes for these triples. The simulated samples are separate from the cooperative's original sample.

\begin{center}
\begin{tabular}{@{}lr@{}}
    \toprule
    Outcome for the intervals from one simulated sample & Number of samples \\
    \midrule
    All three cover the simulated mean & 178 \\
    Only the 95\% and 99\% intervals cover & 12 \\
    Only the 99\% interval covers & 7 \\
    None of the three covers & 3 \\
    \bottomrule
\end{tabular}
\end{center}

\QuestionsTasks[sec:general-worked]{
    {2}/{sol:c2-compare-three-a}/{\textbf{Task 1 -- Organize the information and calculate SE.} Use the four-column information table and classification convention from the practice procedure; distinguish shared values from level-specific settings.},
    {2}/{sol:c2-compare-three-b}/{\textbf{Task 2 -- Calculate MoE and state the intervals.} Show substitutions, endpoints, widths, and units for all three intervals; interpret them and explain how confidence changes width and centre.},
    {2}/{sol:c2-compare-three-c}/{\textbf{Task 3 -- Calculate observed coverage and answer the multiple-choice question.} Show covering counts, denominators, and rates for each level. Which option gives the rates in the order 90\%, 95\%, 99\% and correctly interprets increasing confidence?}
}

\begin{answeroptions}
    \item 90\%, 95\%, 99\%. Every finite simulation must match its nominal levels exactly.
    \item 89\%, 95\%, 98.5\%. Higher confidence gives wider intervals with the same sample mean as centre; confidence describes long-run coverage of the fixed population mean.
    \item 89\%, 95\%, 98.5\%. Higher confidence moves the interval centre closer to the unknown population mean.
    \item 89\%, 95\%, 98.5\%. The cooperative's calculated 99\% interval has a 99\% probability of containing its fixed population mean.
\end{answeroptions}

\clearpage
\subsubsection*{Task 1 solution. Organize the information and calculate SE}
\label{sol:c2-compare-three-a}

\InformationTable{
    $\CtwoMean{\mu}$ & \CtwoMean{Population mean} & Population parameters & The fixed, unknown mean time for all appliance repairs completed by the cooperative in the specified month, in \CtwoMean{minutes}. The separate simulation uses a known mean of \CtwoMean{51 minutes}. \\
    $\CtwoPoint{\bar{x}}$ & \CtwoPoint{Sample mean} & Sample statistics & $\CtwoPoint{\bar{x}}=\CtwoPoint{50\,\text{minutes}}$, shared by all three intervals. The point estimate of the cooperative's unknown mean repair time. \\
    $\CtwoSize{n}$ & \CtwoSize{Sample size} & Sample statistics & $\CtwoSize{n}=\CtwoSize{64\,\text{repairs}}$, shared. All three intervals use the same original sample. \\
    $\CtwoSD{\sigma}$ & \CtwoSD{Population} \CtwoSD{standard} \CtwoSD{deviation} & Population parameters & $\CtwoSD{\sigma}=\CtwoSD{16\,\text{minutes}}$, shared. Known population parameter. \\
    $\CtwoConfidence{C}$ & \CtwoConfidence{Confidence level} & Procedure Settings & \CtwoConfidence{90\%}: $\CtwoConfidence{C}=\CtwoConfidence{0.90}$.\par \CtwoConfidence{95\%}: $\CtwoConfidence{C}=\CtwoConfidence{0.95}$.\par \CtwoConfidence{99\%}: $\CtwoConfidence{C}=\CtwoConfidence{0.99}$. \\
    $\CtwoSignificance{\alpha}$ & \CtwoSignificance{Significance level} & Procedure Settings & \CtwoConfidence{90\%}: $\CtwoSignificance{\alpha}=1-\CtwoConfidence{0.90}=\CtwoSignificance{0.10}$.\par \CtwoConfidence{95\%}: $\CtwoSignificance{\alpha}=1-\CtwoConfidence{0.95}=\CtwoSignificance{0.05}$.\par \CtwoConfidence{99\%}: $\CtwoSignificance{\alpha}=1-\CtwoConfidence{0.99}=\CtwoSignificance{0.01}$. Total area outside the central region. \\
    $\CtwoSignificance{\alpha}/2$ & \CtwoSignificance{Tail area} & Procedure Settings & \CtwoConfidence{90\%}: $\CtwoSignificance{\alpha}/2=\CtwoSignificance{0.05}$.\par \CtwoConfidence{95\%}: $\CtwoSignificance{\alpha}/2=\CtwoSignificance{0.025}$.\par \CtwoConfidence{99\%}: $\CtwoSignificance{\alpha}/2=\CtwoSignificance{0.005}$. Area in each tail. \\
    $\CtwoCritical{z^*}$ & \CtwoCritical{Critical value} & Procedure Settings & \CtwoConfidence{90\%}: $\CtwoCritical{z^*}=\CtwoCritical{1.645}$.\par \CtwoConfidence{95\%}: $\CtwoCritical{z^*}=\CtwoCritical{1.960}$.\par \CtwoConfidence{99\%}: $\CtwoCritical{z^*}=\CtwoCritical{2.576}$. Standard normal multipliers. \\
    $\CtwoSE{\mathrm{SE}}$ & \CtwoSE{Standard error} & Sample statistics & $\CtwoSE{\mathrm{SE}}=\CtwoSD{\sigma}/\sqrt{\CtwoSize{n}}=\CtwoSD{16}/\sqrt{\CtwoSize{64}}=\CtwoSE{2\,\text{minutes}}$, shared. The standard deviation of sample means over repeated samples of this size. \\
}

\par\vspace{3mm}\secInicio[prob:c2-compare-three][sec:general-worked]

\clearpage
\subsubsection*{Task 2 solution. Calculate MoE and state the intervals}
\label{sol:c2-compare-three-b}

For each level, $\CtwoMoE{\mathrm{MoE}}=\CtwoCritical{z^*}\times\CtwoSE{\mathrm{SE}}$, the endpoints are $\CtwoPoint{\bar{x}}\pm\CtwoMoE{\mathrm{MoE}}$, and the \CtwoWidth{full width} is $2\CtwoMoE{\mathrm{MoE}}$.

\EstimationTable{
    \CtwoConfidence{90\%} interval & $\CtwoMoE{\mathrm{MoE}}=\CtwoCritical{1.645}(\CtwoSE{2})=\CtwoMoE{3.29\,\text{minutes}}$.\par $[\CtwoPoint{50}-\CtwoMoE{3.29},\CtwoPoint{50}+\CtwoMoE{3.29}]=[\CtwoEndpoint{46.71},\CtwoEndpoint{53.29}]$ \CtwoEndpoint{minutes}.\par \CtwoWidth{Width}: $2(\CtwoMoE{3.29})=\CtwoWidth{6.58\,\text{minutes}}$. & The narrowest range of plausible values for the cooperative's mean repair time under the model. \\
    \CtwoConfidence{95\%} interval & $\CtwoMoE{\mathrm{MoE}}=\CtwoCritical{1.960}(\CtwoSE{2})=\CtwoMoE{3.92\,\text{minutes}}$.\par $[\CtwoPoint{50}-\CtwoMoE{3.92},\CtwoPoint{50}+\CtwoMoE{3.92}]=[\CtwoEndpoint{46.08},\CtwoEndpoint{53.92}]$ \CtwoEndpoint{minutes}.\par \CtwoWidth{Width}: $2(\CtwoMoE{3.92})=\CtwoWidth{7.84\,\text{minutes}}$. & A wider range of plausible values for the same unknown mean under the model. \\
    \CtwoConfidence{99\%} interval & $\CtwoMoE{\mathrm{MoE}}=\CtwoCritical{2.576}(\CtwoSE{2})=\CtwoMoE{5.152\,\text{minutes}}$.\par $[\CtwoPoint{50}-\CtwoMoE{5.152},\CtwoPoint{50}+\CtwoMoE{5.152}]=[\CtwoEndpoint{44.848},\CtwoEndpoint{55.152}]$ \CtwoEndpoint{minutes}.\par \CtwoWidth{Width}: $2(\CtwoMoE{5.152})=\CtwoWidth{10.304\,\text{minutes}}$. & The widest range of plausible values for the same unknown mean under the model. \\
}

\begin{center}
\begin{tikzpicture}
\begin{axis}[
    width=0.88\linewidth,height=4.7cm,axis lines=left,
    xmin=43,xmax=57,ymin=0.35,ymax=3.7,
    xlabel={Mean appliance repair time (minutes)},
    ytick={1,2,3},
    yticklabels={\CtwoConfidence{90\%},\CtwoConfidence{95\%},\CtwoConfidence{99\%}},
    xtick={44,46,48,50,52,54,56},tick label style={font=\small},label style={font=\small}]
    \addplot[black,dashed] coordinates {(50,0.35) (50,3.3)};
    \node[anchor=south,font=\small] at (axis cs:50,3.3) {$\CtwoPoint{\bar{x}}$};
    \addplot[black,thick,mark=|,mark size=5pt] coordinates {(46.71,1) (53.29,1)};
    \addplot[black,only marks,mark=*,mark size=2pt] coordinates {(50,1)};
    \node[anchor=south,font=\small] at (axis cs:46.71,1.05) {$\CtwoEndpoint{46.71}$};
    \node[anchor=south,font=\small] at (axis cs:53.29,1.05) {$\CtwoEndpoint{53.29}$};
    \addplot[black,thick,mark=|,mark size=5pt] coordinates {(46.08,2) (53.92,2)};
    \addplot[black,only marks,mark=*,mark size=2pt] coordinates {(50,2)};
    \node[anchor=south,font=\small] at (axis cs:46.08,2.05) {$\CtwoEndpoint{46.08}$};
    \node[anchor=south,font=\small] at (axis cs:53.92,2.05) {$\CtwoEndpoint{53.92}$};
    \addplot[black,thick,mark=|,mark size=5pt] coordinates {(44.848,3) (55.152,3)};
    \addplot[black,only marks,mark=*,mark size=2pt] coordinates {(50,3)};
    \node[anchor=south,font=\small] at (axis cs:44.848,3.05) {$\CtwoEndpoint{44.848}$};
    \node[anchor=south,font=\small] at (axis cs:55.152,3.05) {$\CtwoEndpoint{55.152}$};
\end{axis}
\end{tikzpicture}
\end{center}


These are three intervals from the same sample. Their centre markers and dashed line all mark the \CtwoPoint{sample mean} $\CtwoPoint{\bar{x}}=\CtwoPoint{50\,\text{minutes}}$; the cooperative's \CtwoMean{population mean} $\CtwoMean{\mu}$ remains unknown. With the same \CtwoSE{standard error}, increasing confidence increases $\CtwoCritical{z^*}$, hence the \CtwoMoE{margin of error} and \CtwoWidth{width}, while the centre stays fixed. The intervals are nested: the \CtwoConfidence{90\%} interval lies within the \CtwoConfidence{95\%} interval, which lies within the \CtwoConfidence{99\%} interval.

\par\vspace{3mm}\secInicio[prob:c2-compare-three][sec:general-worked]

\clearpage
\noindent\begin{minipage}{\linewidth}
\subsubsection*{Task 3 solution. Calculate coverage and interpret confidence}
\label{sol:c2-compare-three-c}

\textbf{Correct option B.} The mutually exclusive counts total $178+12+7+3=200$ simulated samples. Each level has its own denominator of 200 intervals:
\begin{align*}
    \CtwoConfidence{90\%}:\quad &\frac{178}{200}\times100\%=\CtwoCoverage{89\%},\\
    \CtwoConfidence{95\%}:\quad &\frac{178+12}{200}\times100\%=\frac{190}{200}\times100\%=\CtwoCoverage{95\%},\\
    \CtwoConfidence{99\%}:\quad &\frac{178+12+7}{200}\times100\%=\frac{197}{200}\times100\%=\CtwoCoverage{98.5\%}.
\end{align*}
The evidence is consistent with nested intervals: whenever a narrower interval covers the simulated mean, its wider partners also cover it. Within each simulated triple, the centre is that fresh sample's mean; it can vary between samples.

The \CtwoCoverage{observed rates} describe this finite run and need not match the \CtwoConfidence{nominal levels} exactly. Under the model, about \CtwoConfidence{90\%}, \CtwoConfidence{95\%}, and \CtwoConfidence{99\%} of the respective intervals cover the fixed population mean over many repetitions. Higher confidence requires a \CtwoWidth{wider interval} from a given sample; it does not move the \CtwoPoint{sample mean}. Once calculated, each interval either \CtwoCoverage{contains the fixed mean} or \CtwoCoverage{misses it}; a confidence level is not a probability assigned to that particular fixed interval containing the mean.

\par\vspace{3mm}\secInicio[prob:c2-compare-three][sec:general-worked]
\end{minipage}\par

\clearpage

\section{Advanced Problems}
\label{sec:advanced}

\SectionProblemsTableOfContents{
    \SectionProblemLink{prob:c2-4}{Paired intervals for delivery times} \\
    \SectionProblemLink{prob:c2-5}{Sample size and produce spending}
}

\subsection{Paired intervals for delivery times}
\label{prob:c2-4}

A cooperative takes an independent random sample of $n=36$ orders delivered in October. Their sample mean is $\bar{x}=48$ minutes. Assume normal delivery times with known population standard deviation $\sigma=12$ minutes. Both intervals use the same sample; each nominal level $C$ sets its own $\alpha=1-C$. The standard normal critical value is $z^*=1.645$ for 90\% confidence and $z^*=1.960$ for 95\% confidence.

In a simulation with a known fixed mean $\mu=48$ minutes and the same standard deviation, the cooperative computes a 90\% and a 95\% interval from each of 100 fresh samples of size $n=36$. The row and column labels give the nominal levels. The cells count the observed outcomes for the paired intervals.

\begin{center}
\begin{tabular}{@{}lrr@{}}
    \toprule
    & 95\% interval covers & 95\% interval misses \\
    \midrule
    90\% interval covers & 89 & 0 \\
    90\% interval misses & 7 & 4 \\
    \bottomrule
\end{tabular}
\end{center}

\QuestionsTasks[sec:advanced]{
    {2}/{sol:c2-4-a}/{\textbf{Task 1 -- Organize the information and calculate SE.}},
    {2}/{sol:c2-4-b}/{\textbf{Task 2 -- Calculate MoE and state the interval.}},
    {2}/{sol:c2-4-c}/{\textbf{Task 3 --} Which option gives the observed coverage rates for the 90\% and 95\% procedures, in that order, and correctly uses the way each pair is constructed to judge whether a 90\% interval could cover the mean while its 95\% partner misses?}
}

\begin{answeroptions}
    \item 90\% and 95\%. The 95\% partner must cover whenever the 90\% interval covers because it is wider with the same centre.
    \item 89\% and 96\%. The 95\% partner can miss even when the 90\% interval covers because some 95\% intervals miss the mean.
    \item 89\% and 96\%. The 95\% partner must cover whenever the 90\% interval covers because it is wider with the same centre.
    \item 89\% and 89\%. The 95\% partner must cover whenever the 90\% interval covers because it is wider with the same centre.
\end{answeroptions}

\clearpage
\subsubsection*{Task 1 solution. Organize the information and calculate SE}
\label{sol:c2-4-a}

\InformationTable{
    $\CtwoMean{\mu}$ & \CtwoMean{Population mean} & Population parameters & The mean delivery time of all the cooperative's orders delivered in October, in \CtwoMean{minutes}. The fixed, unknown mean estimated by both intervals. \\
    $\CtwoPoint{\bar{x}}$ & \CtwoPoint{Sample mean} & Sample statistics & $\CtwoPoint{\bar{x}}=\CtwoPoint{48\,\text{minutes}}$. A sample statistic giving a single estimate of the mean October delivery time. \\
    $\CtwoSize{n}$ & \CtwoSize{Sample size} & Sample statistics & $\CtwoSize{n}=\CtwoSize{36\,\text{orders}}$, shared. Both intervals use the same sample. \\
    $\CtwoSD{\sigma}$ & \CtwoSD{Population} \CtwoSD{standard} \CtwoSD{deviation} & Population parameters & $\CtwoSD{\sigma}=\CtwoSD{12\,\text{minutes}}$, shared. Known population parameter. \\
    $\CtwoConfidence{C}$ & \CtwoConfidence{Confidence level} & Procedure Settings & \CtwoConfidence{90\%} procedure: $\CtwoConfidence{C}=\CtwoConfidence{0.90}$.\par \CtwoConfidence{95\%} procedure: $\CtwoConfidence{C}=\CtwoConfidence{0.95}$. Nominal confidence level. \\
    $\CtwoSignificance{\alpha}$ & \CtwoSignificance{Significance level} & Procedure Settings & \CtwoConfidence{90\%} procedure: $\CtwoSignificance{\alpha}=1-\CtwoConfidence{0.90}=\CtwoSignificance{0.10}$.\par \CtwoConfidence{95\%} procedure: $\CtwoSignificance{\alpha}=1-\CtwoConfidence{0.95}=\CtwoSignificance{0.05}$. Total area outside the central region. \\
    $\CtwoSignificance{\alpha}/2$ & \CtwoSignificance{Tail area} & Procedure Settings & \CtwoConfidence{90\%} procedure: $\CtwoSignificance{\alpha}/2=\CtwoSignificance{0.05}$ in each tail.\par \CtwoConfidence{95\%} procedure: $\CtwoSignificance{\alpha}/2=\CtwoSignificance{0.025}$ in each tail. Area in each tail of the two-sided procedure. \\
    $\CtwoCritical{z^*}$ & \CtwoCritical{Critical value} & Procedure Settings & \CtwoConfidence{90\%} procedure: $\CtwoCritical{z^*}=\CtwoCritical{1.645}$.\par \CtwoConfidence{95\%} procedure: $\CtwoCritical{z^*}=\CtwoCritical{1.960}$. Corresponding standard normal critical value. \\
    $\CtwoSE{\mathrm{SE}}$ & \CtwoSE{Standard error} & Sample statistics & $\CtwoSE{\mathrm{SE}}=\CtwoSD{\sigma}/\sqrt{\CtwoSize{n}}=\CtwoSD{12}/\sqrt{\CtwoSize{36}}=\CtwoSE{2\,\text{minutes}}$. Both intervals use the same sample and \CtwoSE{standard error}. \\
}

\par\vspace{3mm}\secInicio[prob:c2-4][sec:advanced]

\clearpage
\subsubsection*{Task 2 solution. Calculate MoE and state the interval}
\label{sol:c2-4-b}

For each procedure, $\CtwoMoE{\mathrm{MoE}}=\CtwoCritical{z^*}\times\CtwoSE{\mathrm{SE}}$ and the interval is $[\CtwoPoint{\bar{x}}-\CtwoMoE{\mathrm{MoE}},\CtwoPoint{\bar{x}}+\CtwoMoE{\mathrm{MoE}}]$.

\EstimationTable{
    \CtwoConfidence{90\%} interval & $\CtwoMoE{\mathrm{MoE}}=\CtwoCritical{z^*}\CtwoSE{\mathrm{SE}}=\CtwoCritical{1.645}(\CtwoSE{2})=\CtwoMoE{3.29\,\text{minutes}}$. The \CtwoEndpoint{endpoints} from $\CtwoPoint{48}\pm\CtwoMoE{3.29}$ give $[\CtwoEndpoint{44.71},\CtwoEndpoint{51.29}]$ \CtwoEndpoint{minutes}. & The narrower range of plausible values for the mean October delivery time under the model. \\
    \CtwoConfidence{95\%} interval & $\CtwoMoE{\mathrm{MoE}}=\CtwoCritical{z^*}\CtwoSE{\mathrm{SE}}=\CtwoCritical{1.960}(\CtwoSE{2})=\CtwoMoE{3.92\,\text{minutes}}$. The \CtwoEndpoint{endpoints} from $\CtwoPoint{48}\pm\CtwoMoE{3.92}$ give $[\CtwoEndpoint{44.08},\CtwoEndpoint{51.92}]$ \CtwoEndpoint{minutes}. & A wider range of plausible values for the same mean under the model, extending farther on both sides of \CtwoPoint{48}. \\
}

The widths are $2(\CtwoMoE{3.29})=\CtwoWidth{6.58\,\text{minutes}}$ and $2(\CtwoMoE{3.92})=\CtwoWidth{7.84\,\text{minutes}}$. At the same sample size and SE, the higher confidence level gives a wider interval.

\par\vspace{3mm}\secInicio[prob:c2-4][sec:advanced]


\noindent\begin{minipage}{\linewidth}
\subsubsection*{Task 3 solution. Calculate coverage and interpret confidence}
\label{sol:c2-4-c}

\textbf{Correct option C.} The \CtwoCoverage{observed coverage} for the \CtwoConfidence{90\%} procedure is $(89+0)/100\times100\%=\CtwoCoverage{89\%}$ from the covering row. For the \CtwoConfidence{95\%} procedure it is $(89+7)/100\times100\%=\CtwoCoverage{96\%}$ from the covering column.

Within each pair, the intervals share a centre and \CtwoSE{standard error}. The \CtwoCritical{larger critical value} for \CtwoConfidence{95\%} gives a larger $\CtwoMoE{\mathrm{MoE}}$, so the \CtwoConfidence{95\%} interval contains the entire \CtwoConfidence{90\%} interval. The \CtwoConfidence{95\%} partner therefore covers whenever the \CtwoConfidence{90\%} interval covers, which accounts for the zero in the table.

Under the model assumptions, about \CtwoConfidence{90\%} and \CtwoConfidence{95\%} of the respective intervals would cover the fixed mean over repeated sampling. Each calculated interval either \CtwoCoverage{contains the mean} or \CtwoCoverage{misses it}.

\par\vspace{3mm}\secInicio[prob:c2-4][sec:advanced]
\end{minipage}\par

\clearpage

\subsection{Sample size and produce spending}
\label{prob:c2-5}

A municipal survey estimates the mean amount households in one district spent on fresh produce during a specified week. Two teams take separate random samples. Assume independent household observations and normally distributed spending with known population standard deviation $\sigma=50$ thousand COP. Both teams use the same nominal level $C=0.95$ and significance level $\alpha=1-C$, with the standard normal critical value $z^*=1.960$ for 95\% confidence.

\begin{center}
\begin{tabular}{@{}lrr@{}}
    \toprule
    Team & Sample size $n$ & Mean $\bar{x}$ (thousand COP) \\
    \midrule
    A & 25 & 205 \\
    B & 100 & 205 \\
    \bottomrule
\end{tabular}
\end{center}

The teams simulate 200 fresh samples at each size from a population with known fixed mean $\mu=205$ thousand COP and the same standard deviation. Each procedure uses its own 200 intervals as its denominator. Team A records 192 intervals that cover the mean. Team B records 187.

\QuestionsTasks[sec:advanced]{
    {2}/{sol:c2-5-a}/{\textbf{Task 1 -- Organize the information and calculate SE.}},
    {2}/{sol:c2-5-b}/{\textbf{Task 2 -- Calculate MoE and state the interval.}},
    {2}/{sol:c2-5-c}/{\textbf{Task 3 --} Which option correctly compares Team B's interval width with Team A's, gives the observed coverage rates for A and B, and assesses the claim that Team A's higher observed rate shows better long-run coverage?}
}

\begin{answeroptions}
    \item B's interval is half as wide. A covers 96\% and B covers 93.5\% in the simulations, so A has better long-run coverage.
    \item B's interval is half as wide. A covers 96\% and B covers 93.5\% in the simulations. Both procedures cover about 95\% of the time in the long run under the model.
    \item B's interval is one-quarter as wide. A covers 96\% and B covers 93.5\% in the simulations. Both procedures cover about 95\% of the time in the long run under the model.
    \item B's interval is half as wide. A and B each cover 95\% in the simulations and cover about 95\% of the time in the long run under the model.
\end{answeroptions}

\clearpage
\subsubsection*{Task 1 solution. Organize the information and calculate SE}
\label{sol:c2-5-a}

\InformationTable{
    $\CtwoMean{\mu}$ & \CtwoMean{Population mean} & Population parameters & The mean spending on fresh produce among all households in the district during the specified week, in \CtwoMean{thousand COP}. The fixed, unknown district mean both teams want to estimate. \\
    $\CtwoPoint{\bar{x}}$ & \CtwoPoint{Sample mean} & Sample statistics & Team A and Team B both give $\CtwoPoint{\bar{x}}=\CtwoPoint{205\,\text{thousand COP}}$. Sample statistics giving the same single estimate of the mean weekly produce spending from both samples. \\
    $\CtwoSize{n}$ & \CtwoSize{Sample size} & Sample statistics & Team A: $\CtwoSize{n}=\CtwoSize{25}$; Team B: $\CtwoSize{n}=\CtwoSize{100\,\text{households}}$. Separate samples of different sizes. \\
    $\CtwoSD{\sigma}$ & \CtwoSD{Population} \CtwoSD{standard} \CtwoSD{deviation} & Population parameters & $\CtwoSD{\sigma}=\CtwoSD{50\,\text{thousand COP}}$, shared. Known population parameter. \\
    $\CtwoConfidence{C}$ & \CtwoConfidence{Confidence level} & Procedure Settings & Team A and Team B (\CtwoConfidence{95\%}): $\CtwoConfidence{C}=\CtwoConfidence{0.95}$. Nominal confidence level. \\
    $\CtwoSignificance{\alpha}$ & \CtwoSignificance{Significance level} & Procedure Settings & Team A and Team B (\CtwoConfidence{95\%}): $\CtwoSignificance{\alpha}=1-\CtwoConfidence{0.95}=\CtwoSignificance{0.05}$. Total area outside the central region. \\
    $\CtwoSignificance{\alpha}/2$ & \CtwoSignificance{Tail area} & Procedure Settings & Team A and Team B (\CtwoConfidence{95\%}): $\CtwoSignificance{\alpha}/2=\CtwoSignificance{0.025}$ in each tail. Area in each tail of the two-sided procedure. \\
    $\CtwoCritical{z^*}$ & \CtwoCritical{Critical value} & Procedure Settings & Team A and Team B (\CtwoConfidence{95\%}): $\CtwoCritical{z^*}=\CtwoCritical{1.960}$. Corresponding standard normal critical value. \\
    $\CtwoSE{\mathrm{SE}}$ & \CtwoSE{Standard error} & Sample statistics & Using $\CtwoSE{\mathrm{SE}}=\CtwoSD{\sigma}/\sqrt{\CtwoSize{n}}$, Team A gets $\CtwoSD{50}/\sqrt{\CtwoSize{25}}=\CtwoSE{10}$ and Team B gets $\CtwoSD{50}/\sqrt{\CtwoSize{100}}=\CtwoSE{5}$, both in \CtwoSE{thousand COP}. Team B's \CtwoPoint{sample mean} varies less across repeated samples. \\
}

\par\vspace{3mm}\secInicio[prob:c2-5][sec:advanced]

\clearpage
\subsubsection*{Task 2 solution. Calculate MoE and state the interval}
\label{sol:c2-5-b}

For each procedure, $\CtwoMoE{\mathrm{MoE}}=\CtwoCritical{z^*}\times\CtwoSE{\mathrm{SE}}$ and the interval is $[\CtwoPoint{\bar{x}}-\CtwoMoE{\mathrm{MoE}},\CtwoPoint{\bar{x}}+\CtwoMoE{\mathrm{MoE}}]$.

\EstimationTable{
    \CtwoConfidence{95\%} interval (Team A) & $\CtwoMoE{\mathrm{MoE}}=\CtwoCritical{z^*}\CtwoSE{\mathrm{SE}}=\CtwoCritical{1.960}(\CtwoSE{10})=\CtwoMoE{19.6\,\text{thousand COP}}$. The \CtwoEndpoint{endpoints} from $\CtwoPoint{205}\pm\CtwoMoE{19.6}$ give $[\CtwoEndpoint{185.4},\CtwoEndpoint{224.6}]$ \CtwoEndpoint{thousand COP}. & The wider range of plausible values for the district mean under the model. \\
    \CtwoConfidence{95\%} interval (Team B) & $\CtwoMoE{\mathrm{MoE}}=\CtwoCritical{z^*}\CtwoSE{\mathrm{SE}}=\CtwoCritical{1.960}(\CtwoSE{5})=\CtwoMoE{9.8\,\text{thousand COP}}$. The \CtwoEndpoint{endpoints} from $\CtwoPoint{205}\pm\CtwoMoE{9.8}$ give $[\CtwoEndpoint{195.2},\CtwoEndpoint{214.8}]$ \CtwoEndpoint{thousand COP}. & The narrower range of plausible values for the same district mean under the model. \\
}

The \CtwoWidth{widths} are $2(\CtwoMoE{19.6})=\CtwoWidth{39.2}$ and $2(\CtwoMoE{9.8})=\CtwoWidth{19.6\,\text{thousand COP}}$. Team B's \CtwoSize{fourfold sample size} halves its \CtwoSE{standard error} and therefore halves its $\CtwoMoE{\mathrm{MoE}}$ and \CtwoWidth{interval width}.

\par\vspace{3mm}\secInicio[prob:c2-5][sec:advanced]


\noindent\begin{minipage}{\linewidth}
\subsubsection*{Task 3 solution. Calculate coverage and interpret confidence}
\label{sol:c2-5-c}

\textbf{Correct option B.}

The \CtwoCoverage{observed rates} are $(192/200)\times100\%=\CtwoCoverage{96\%}$ for A and $(187/200)\times100\%=\CtwoCoverage{93.5\%}$ for B. These rates can fluctuate across runs. Both procedures use the nominal level of \CtwoConfidence{95\%} and cover when the distance from $\CtwoPoint{\bar{x}}$ to $\CtwoMean{\mu}$ is at most $\CtwoMoE{\mathrm{MoE}}=\CtwoCritical{1.960}\,\CtwoSE{\mathrm{SE}}$. Under the model assumptions, that rule gives both procedures about \CtwoConfidence{95\%} coverage in the long run. Each calculated interval either \CtwoCoverage{contains the mean} or \CtwoCoverage{misses it}.

\par\vspace{3mm}\secInicio[prob:c2-5][sec:advanced]
\end{minipage}\par

\end{document}
